% =====================================================================
%  RESULTADOS E DISCUSSÃO
%  Importe com % =====================================================================
%  RESULTADOS E DISCUSSÃO
%  Importe com % =====================================================================
%  RESULTADOS E DISCUSSÃO
%  Importe com % =====================================================================
%  RESULTADOS E DISCUSSÃO
%  Importe com \input{resultados} a partir do seu .tex principal.
%
%  Este arquivo NÃO carrega pacote nenhum e NÃO declara \documentclass.
%  Ele usa apenas o que o protocolo-experimental.tex já exige:
%    booktabs, multirow, graphicx  (e, se você usar as figuras, um
%    \graphicspath{{figuras/}} no preâmbulo ou o caminho corrigido nos
%    \includegraphics abaixo).
%
%  Se a sua classe for `article`, troque \chapter -> \section e rebaixe
%  os demais níveis.
% =====================================================================

\chapter{Resultados e Discussão}
\label{ch:results}

Este capítulo reporta a grade completa descrita no
Capítulo~\ref{ch:protocol}: 162 treinos (3 arquiteturas $\times$ 2 conjuntos
$\times$ 3 frações de dados $\times$ 3 regimes de augmentação $\times$ 3
sementes), cada um avaliado sob três condições de teste --- limpo, rotação
grande e composição de todas as transformações não vistas --- com e sem
augmentação forte aplicada em teste. São 324 células de resultado e 972
avaliações individuais por semente.

A leitura que segue está organizada em torno de uma decisão metodológica que
precisa ser explicitada antes dos números: \emph{a razão de retenção
$R_c$ da Equação~\eqref{eq:retention} não é usada como métrica principal}.
A Seção~\ref{sec:res:confound} mostra que, nesta grade, $R_c$ é
sistematicamente anticorrelacionada com a acurácia limpa e portanto premia
modelos fracos. As conclusões substantivas repousam sobre a acurácia absoluta
sob deslocamento, lida contra o acaso e contra a linha de base da classe
majoritária.

% ---------------------------------------------------------------------
\section{Desempenho em distribuição}
\label{sec:res:clean}

A Tabela~\ref{tab:res-clean} traz a acurácia de teste limpa. O ordenamento é
estável em toda a grade e não depende da fração de dados.

\begin{table}[t]
  \centering
  \caption{Acurácia de teste limpa (\%) por arquitetura, fração de dados e
  regime de augmentação. Média $\pm$ desvio padrão sobre três sementes.}
  \label{tab:res-clean}
  \small
  \begin{tabular}{llccc}
    \toprule
    \multirow{2}{*}{Modelo} & \multirow{2}{*}{Aug.} & \multicolumn{3}{c}{Fração de dados $f$} \\
    \cmidrule(lr){3-5}
     & & $0{,}33$ & $0{,}66$ & $1{,}00$ \\
    \midrule
    \multicolumn{5}{l}{\textit{CIFAR-10}} \\
    \multirow{3}{*}{Efficient-CapsNet}
      & nenhuma & 63,8 $\pm$ 0,1 & 70,3 $\pm$ 0,4 & 74,2 $\pm$ 0,1 \\
      & padrão  & 78,0 $\pm$ 0,5 & 82,2 $\pm$ 0,1 & 84,8 $\pm$ 0,4 \\
      & forte   & 76,6 $\pm$ 0,3 & 80,7 $\pm$ 0,1 & 83,0 $\pm$ 0,3 \\
    \multirow{3}{*}{ResNet-18}
      & nenhuma & 66,2 $\pm$ 1,9 & 75,3 $\pm$ 0,9 & 78,8 $\pm$ 1,5 \\
      & padrão  & 83,3 $\pm$ 0,5 & 88,6 $\pm$ 0,4 & 90,8 $\pm$ 0,1 \\
      & forte   & \textbf{86,8 $\pm$ 0,6} & \textbf{91,5 $\pm$ 0,1} & \textbf{93,3 $\pm$ 0,1} \\
    \multirow{3}{*}{DeiT-Tiny}
      & nenhuma & 52,8 $\pm$ 0,7 & 58,8 $\pm$ 0,5 & 60,9 $\pm$ 2,2 \\
      & padrão  & 59,2 $\pm$ 1,3 & 61,6 $\pm$ 6,6 & 70,2 $\pm$ 3,5 \\
      & forte   & 49,6 $\pm$ 4,0 & 57,8 $\pm$ 8,0 & 68,7 $\pm$ 10,4 \\
    \midrule
    \multicolumn{5}{l}{\textit{D-Fire}} \\
    \multirow{3}{*}{Efficient-CapsNet}
      & nenhuma & 87,3 $\pm$ 0,5 & 89,2 $\pm$ 0,2 & 90,2 $\pm$ 0,5 \\
      & padrão  & \textbf{89,0 $\pm$ 0,5} & 92,8 $\pm$ 0,4 & 93,6 $\pm$ 0,4 \\
      & forte   & 86,4 $\pm$ 0,0 & 90,3 $\pm$ 0,4 & 91,9 $\pm$ 0,3 \\
    \multirow{3}{*}{ResNet-18}
      & nenhuma & 84,4 $\pm$ 1,7 & 88,7 $\pm$ 0,6 & 91,2 $\pm$ 0,4 \\
      & padrão  & 88,8 $\pm$ 0,9 & 92,7 $\pm$ 0,3 & 94,1 $\pm$ 0,5 \\
      & forte   & \textbf{90,0 $\pm$ 0,7} & \textbf{93,8 $\pm$ 0,2} & \textbf{94,7 $\pm$ 0,2} \\
    \multirow{3}{*}{DeiT-Tiny}
      & nenhuma & 82,2 $\pm$ 0,7 & 84,0 $\pm$ 2,5 & 84,6 $\pm$ 2,9 \\
      & padrão  & 76,0 $\pm$ 7,4 & 78,9 $\pm$ 8,8 & 81,0 $\pm$ 5,4 \\
      & forte   & 73,9 $\pm$ 2,0 & 75,1 $\pm$ 5,3 & 77,5 $\pm$ 8,1 \\
    \bottomrule
  \end{tabular}
\end{table}

No CIFAR-10 a ResNet-18 vence em todas as nove combinações de fração e
augmentação, com margem entre $4{,}4$ e $10{,}5$ pontos percentuais (pp) sobre o
Efficient-CapsNet. O melhor resultado capsular ($84{,}8\%$, augmentação padrão,
$f = 1{,}00$) fica $8{,}5$~pp abaixo do melhor resultado convolucional
($93{,}3\%$). O DeiT-Tiny é o pior dos três por larga margem, o que é esperado:
transformadores de visão treinados do zero, sem pré-treino e com 45.000 imagens
de $32\times32$, operam muito abaixo do seu regime.

No D-Fire o quadro é diferente. A tarefa é binária e mais fácil, e o
Efficient-CapsNet fica a $1{,}1$~pp da ResNet-18 na melhor configuração de cada
um ($93{,}6\%$ contra $94{,}7\%$). Em duas células --- augmentação nenhuma com
$f = 0{,}33$ e $f = 0{,}66$, e augmentação padrão com $f = 0{,}33$ --- o modelo
capsular tem a maior média, mas nenhuma dessas vantagens excede o dobro do
ruído entre sementes.

Duas observações estruturais valem o registro. A primeira: a ResNet-18 é a
única arquitetura que melhora ao subir de augmentação padrão para forte.
Em $f = 1{,}00$, a passagem custa $1{,}8$~pp ao Efficient-CapsNet no CIFAR-10 e
$1{,}7$~pp no D-Fire, e custa $1{,}6$ e $3{,}5$~pp ao DeiT-Tiny, enquanto dá
$2{,}5$ e $0{,}6$~pp à ResNet-18. A segunda: o desvio padrão entre sementes do
Efficient-CapsNet é a menor de todas em quase toda a tabela, ponto retomado na
Seção~\ref{sec:res:seeds}.

% ---------------------------------------------------------------------
\section{Robustez sob transformações não vistas}
\label{sec:res:robust}

A Tabela~\ref{tab:res-robust} reporta a condição $f = 1{,}00$ sob as duas
sondas de deslocamento. Além da acurácia absoluta, listamos $R_c$ e a
\emph{habilidade retida}
\begin{equation}
  S_c = \frac{A^{c} - A^{\text{acaso}}}{A^{\text{clean}} - A^{\text{acaso}}},
  \label{eq:skill}
\end{equation}
que desconta o acerto trivial antes de formar a razão. A diferença entre as
duas colunas é instrutiva: $R_c$ dá crédito por acertar ao acaso, $S_c$ não.

\begin{table}[t]
  \centering
  \caption{Desempenho sob transformações não vistas em $f = 1{,}00$, sem
  augmentação em teste. $A^{c}$ é a acurácia absoluta (\%), $R_c$ a razão de
  retenção (\%) e $S_c$ a habilidade retida (\%) da Equação~\eqref{eq:skill}.
  Acaso: $10\%$ no CIFAR-10, $50\%$ no D-Fire.}
  \label{tab:res-robust}
  \small
  \begin{tabular}{llcccccc}
    \toprule
    \multirow{2}{*}{Modelo} & \multirow{2}{*}{Aug.}
      & \multicolumn{3}{c}{Rotação grande} & \multicolumn{3}{c}{Todas as transformações} \\
    \cmidrule(lr){3-5}\cmidrule(lr){6-8}
     & & $A^{c}$ & $R_c$ & $S_c$ & $A^{c}$ & $R_c$ & $S_c$ \\
    \midrule
    \multicolumn{8}{l}{\textit{CIFAR-10}} \\
    \multirow{3}{*}{Efficient-CapsNet}
      & nenhuma & 38,1 & 51,3 & 43,7 & 12,6 & 17,0 & 4,1 \\
      & padrão  & 52,9 & 62,4 & 57,4 & 12,8 & 15,1 & 3,8 \\
      & forte   & 51,4 & 61,9 & 56,6 & 13,8 & 16,7 & 5,2 \\
    \multirow{3}{*}{ResNet-18}
      & nenhuma & 36,4 & 46,2 & 38,4 & 11,6 & 14,7 & 2,3 \\
      & padrão  & 57,8 & 63,6 & 59,1 & 12,3 & 13,6 & 2,9 \\
      & forte   & \textbf{60,0} & 64,4 & 60,1 & 13,4 & 14,4 & 4,1 \\
    \multirow{3}{*}{DeiT-Tiny}
      & nenhuma & 40,0 & 65,8 & 59,0 & \textbf{16,4} & 26,9 & 12,6 \\
      & padrão  & 52,8 & 75,1 & 71,0 & 15,9 & 22,6 & 9,8 \\
      & forte   & 52,0 & \textbf{75,7} & \textbf{71,6} & \textbf{16,8} & \textbf{24,5} & \textbf{11,6} \\
    \midrule
    \multicolumn{8}{l}{\textit{D-Fire}} \\
    \multirow{3}{*}{Efficient-CapsNet}
      & nenhuma & 61,4 & 68,1 & 28,4 & 53,1 & 58,9 & 7,7 \\
      & padrão  & 77,6 & 82,9 & 63,2 & 56,1 & 59,9 & 13,9 \\
      & forte   & 80,0 & 87,0 & 71,5 & 60,5 & 65,8 & 25,0 \\
    \multirow{3}{*}{ResNet-18}
      & nenhuma & 63,4 & 69,5 & 32,4 & 52,5 & 57,5 & 6,0 \\
      & padrão  & 78,7 & 83,7 & 65,2 & 54,9 & 58,4 & 11,2 \\
      & forte   & \textbf{80,9} & 85,4 & 69,1 & \textbf{60,7} & 64,1 & 23,9 \\
    \multirow{3}{*}{DeiT-Tiny}
      & nenhuma & 63,5 & 75,0 & 38,9 & 52,9 & 62,5 & 8,4 \\
      & padrão  & 78,4 & 96,9 & 91,8 & 53,7 & 66,4 & 12,1 \\
      & forte   & 76,0 & \textbf{98,0} & \textbf{94,4} & 59,9 & \textbf{77,2} & \textbf{35,9} \\
    \bottomrule
  \end{tabular}
\end{table}

\subsection{A sonda de rotação}

A rotação grande é o teste direto de H1, e o resultado é desfavorável à
hipótese. No CIFAR-10 com augmentação padrão ou forte, a ResNet-18 tem a maior
acurácia absoluta sob rotação --- $57{,}8\%$ e $60{,}0\%$ contra $52{,}9\%$ e
$51{,}4\%$ do Efficient-CapsNet. A vantagem convolucional sob rotação
($+8{,}6$~pp em $f = 1{,}00$, augmentação forte) é maior, em valor absoluto, do
que qualquer vantagem capsular observada em toda a grade.

Apenas no regime sem augmentação nenhuma o Efficient-CapsNet fica à frente sob
rotação, e por pouco: $38{,}1\%$ contra $36{,}4\%$ da ResNet-18. Em teste
pareado por semente sobre todas as 27 células do CIFAR-10, a diferença de
retenção entre Efficient-CapsNet e ResNet-18 sob rotação é de $+1{,}1$~pp com
$p = 0{,}13$ --- indistinguível de zero. No D-Fire a mesma comparação dá
$-1{,}5$~pp com $p = 0{,}013$, ou seja, uma desvantagem capsular
estatisticamente detectável. Contra o DeiT-Tiny a diferença de retenção é de
$-15{,}1$~pp no CIFAR-10 e $-11{,}7$~pp no D-Fire, ambas com $p < 10^{-12}$.

A Figura~\ref{fig:res-dumbbell} torna o padrão visível: as barras do
Efficient-CapsNet e da ResNet-18 partem de pontos diferentes e caem
aproximadamente na mesma proporção, sem que a distância entre elas se feche.

\begin{figure}[t]
  \centering
  \includegraphics[width=\textwidth]{fig01_shift_dumbbell_cifar_10.png}
  \caption{CIFAR-10: acurácia limpa (círculo aberto) e acurácia sob
  deslocamento (marcador preenchido) para cada combinação de arquitetura,
  fração de dados e regime de augmentação. O comprimento da barra é a
  degradação absoluta.}
  \label{fig:res-dumbbell}
\end{figure}

\subsection{A composição de todas as transformações é um piso, não uma medida}
\label{sec:res:floor}

Sob a composição das seis corrupções, os três modelos deixam de funcionar. No
CIFAR-10 a acurácia cai para a faixa de $11{,}1\%$ a $19{,}1\%$, contra
$10{,}0\%$ de acaso, em todas as 54 células. No D-Fire os valores ficam entre
$-2{,}2$ e $+7{,}8$~pp da linha de base da classe majoritária ($53{,}4\%$); sem
augmentação de treino, 14 das 18 células ficam \emph{abaixo} dela, isto é,
perde para um classificador que ignora a imagem
(Figura~\ref{fig:res-baselines}).

Isso muda o que se pode concluir a partir dessa condição. Comparar retenção
entre modelos que estão todos no nível do acaso compara padrões de colapso, não
robustez. As diferenças de $R_c$ que aparecem ali --- $26{,}9\%$ para o
DeiT-Tiny contra $14{,}7\%$ para a ResNet-18, por exemplo --- correspondem, em
acurácia absoluta, a $16{,}4\%$ contra $11{,}6\%$: quatro pontos e meio acima e
um ponto e meio acima do acaso, respectivamente. Nenhum dos dois é utilizável.

\begin{figure}[t]
  \centering
  \includegraphics[width=\textwidth]{fig11_cross_dataset_vs_baseline.png}
  \caption{Os dois conjuntos contra suas próprias linhas de base triviais.
  Barras abaixo da linha pontilhada perdem para um classificador que ignora a
  imagem. Sob todas as transformações, nenhuma arquitetura se sustenta em
  nenhum dos conjuntos.}
  \label{fig:res-baselines}
\end{figure}

A recomendação que decorre daí é metodológica e vale além deste trabalho:
composições de corrupções só informam enquanto os modelos comparados
permanecerem acima do acaso por margem apreciável. Abaixo disso, a métrica
mede outra coisa. A Seção~\ref{sec:res:collapse} identifica exatamente o que.

\subsection{Por que a retenção não pode ser a métrica principal}
\label{sec:res:confound}

O DeiT-Tiny tem a maior razão de retenção em praticamente toda a
Tabela~\ref{tab:res-robust} e a menor acurácia limpa em toda a
Tabela~\ref{tab:res-clean}. Não é coincidência. A Figura~\ref{fig:res-confound}
mostra a correlação entre acurácia limpa e retenção no CIFAR-10:
$r = -0{,}44$ sob rotação e $r = -0{,}91$ sob todas as transformações. As
inclinações ajustadas são de $-34$ e $-48$~pp de retenção por unidade de
acurácia limpa.

\begin{figure}[t]
  \centering
  \includegraphics[width=\textwidth]{fig03_retention_confound_cifar_10.png}
  \caption{CIFAR-10: razão de retenção contra acurácia limpa. A relação
  negativa significa que a retenção, isolada, recompensa modelos que já eram
  ruins em distribuição.}
  \label{fig:res-confound}
\end{figure}

O mecanismo é aritmético. A acurácia sob deslocamento tem piso no acaso,
enquanto a acurácia limpa não tem teto próximo; um modelo que sai de $60\%$ e
cai para $40\%$ retém $67\%$, e um que sai de $93\%$ e cai para $60\%$ retém
$64\%$ --- embora o segundo esteja $20$~pp à frente onde importa. A
Figura~\ref{fig:res-frontier} mostra a mesma informação no plano
(acurácia limpa, acurácia deslocada): sob rotação a correlação entre os dois
eixos é $+0{,}68$, e sob todas as transformações ela \emph{inverte} de sinal
($-0{,}74$), que é a assinatura do regime de colapso descrito acima.

\begin{figure}[t]
  \centering
  \includegraphics[width=\textwidth]{fig02_robustness_frontier_cifar_10.png}
  \caption{CIFAR-10: acurácia sob deslocamento contra acurácia limpa. A
  diagonal tracejada marca invariância perfeita. Separação vertical entre cores
  na mesma abscissa seria o efeito arquitetural puro; ela é pequena sob rotação
  e some sob todas as transformações.}
  \label{fig:res-frontier}
\end{figure}

Por isso a Tabela~\ref{tab:res-robust} reporta $A^c$ em primeiro lugar, e por
isso H1 é avaliada sobre acurácia absoluta e sobre retenção pareada por
semente, não sobre médias de $R_c$.

% ---------------------------------------------------------------------
\section{Colapso de predição e detecção de erro}
\label{sec:res:collapse}

Dois diagnósticos separam ``degradação graciosa'' de ``falha silenciosa'': a
fração de predições que caem na classe mais predita (\emph{top-pred share}) e a
AUROC da confiança do modelo como detector dos próprios erros. A
Tabela~\ref{tab:res-auroc} traz os dois.

\begin{table}[t]
  \centering
  \caption{AUROC de detecção de erro (a confiança separa acertos de erros?) e
  fração de predições na classe mais predita, médias sobre as três frações de
  dados, sem augmentação em teste. AUROC de $0{,}50$ é o acaso; \emph{share} de
  $0{,}10$ (CIFAR-10) e $0{,}53$ (D-Fire) corresponde a predições
  balanceadas.}
  \label{tab:res-auroc}
  \footnotesize
  \setlength{\tabcolsep}{4pt}
  \begin{tabular}{lccccccccc}
    \toprule
    \multirow{2}{*}{Modelo}
      & \multicolumn{3}{c}{Limpo} & \multicolumn{3}{c}{Rotação grande} & \multicolumn{3}{c}{Todas} \\
    \cmidrule(lr){2-4}\cmidrule(lr){5-7}\cmidrule(lr){8-10}
     & nenh. & padr. & forte & nenh. & padr. & forte & nenh. & padr. & forte \\
    \midrule
    \multicolumn{10}{l}{\textit{CIFAR-10 --- AUROC}} \\
    Efficient-CapsNet & 0,799 & 0,856 & 0,848 & 0,707 & 0,762 & 0,766 & 0,533 & 0,547 & 0,548 \\
    ResNet-18         & 0,788 & 0,851 & 0,867 & 0,697 & 0,790 & 0,804 & 0,540 & 0,527 & 0,540 \\
    DeiT-Tiny         & 0,716 & 0,736 & 0,720 & 0,677 & 0,722 & 0,713 & 0,560 & 0,587 & 0,618 \\
    \addlinespace
    \multicolumn{10}{l}{\textit{CIFAR-10 --- fração na classe mais predita}} \\
    Efficient-CapsNet & 0,12 & 0,11 & 0,11 & 0,24 & 0,18 & 0,20 & 0,39 & 0,48 & 0,43 \\
    ResNet-18         & 0,13 & 0,11 & 0,11 & 0,28 & 0,18 & 0,18 & 0,55 & 0,70 & 0,76 \\
    DeiT-Tiny         & 0,13 & 0,12 & 0,13 & 0,21 & 0,15 & 0,16 & 0,38 & 0,48 & 0,38 \\
    \midrule
    \multicolumn{10}{l}{\textit{D-Fire --- AUROC}} \\
    Efficient-CapsNet & 0,837 & 0,838 & 0,836 & 0,581 & 0,696 & 0,728 & 0,515 & 0,558 & 0,595 \\
    ResNet-18         & 0,826 & 0,845 & 0,881 & 0,610 & 0,736 & 0,745 & 0,532 & 0,562 & 0,601 \\
    DeiT-Tiny         & 0,769 & 0,755 & 0,731 & 0,627 & 0,727 & 0,726 & 0,525 & 0,560 & 0,566 \\
    \addlinespace
    \multicolumn{10}{l}{\textit{D-Fire --- fração na classe mais predita}} \\
    Efficient-CapsNet & 0,56 & 0,54 & 0,55 & 0,78 & 0,64 & 0,62 & 0,75 & 0,60 & 0,58 \\
    ResNet-18         & 0,53 & 0,53 & 0,53 & 0,70 & 0,61 & 0,62 & 0,65 & 0,58 & 0,62 \\
    DeiT-Tiny         & 0,55 & 0,54 & 0,54 & 0,55 & 0,54 & 0,54 & 0,70 & 0,63 & 0,63 \\
    \bottomrule
  \end{tabular}
\end{table}

Em dados limpos, o Efficient-CapsNet tem a melhor detecção de erro do CIFAR-10
sob augmentação padrão ($0{,}856$ contra $0{,}851$ da ResNet-18) e a melhor do
D-Fire no regime sem augmentação ($0{,}837$ contra $0{,}826$), ficando a menos
de um centésimo da ResNet-18 sob augmentação padrão. A norma do vetor de cápsula, portanto,
funciona como escore de confiança pelo menos tão bem quanto o \emph{softmax},
o que é um resultado positivo para o modelo capsular e independente da
acurácia.

Sob deslocamento essa propriedade não sobrevive. Sob rotação a AUROC capsular
cai para $0{,}707$--$0{,}766$ no CIFAR-10 e para $0{,}581$ no D-Fire sem
augmentação --- neste último caso, com $77{,}6\%$ das predições em uma única
classe. Ou seja: o modelo rotacionado no D-Fire responde quase sempre a mesma
coisa e não sabe que está errando. Sob todas as transformações, a AUROC dos
três modelos fica entre $0{,}515$ e $0{,}618$, praticamente o acaso
(Figura~\ref{fig:res-auroc}). A falha, nessa condição, é silenciosa em todas as
arquiteturas.

\begin{figure}[t]
  \centering
  \includegraphics[width=\textwidth]{fig07_error_detection_auroc_cifar_10.png}
  \caption{CIFAR-10: AUROC da confiança como detector dos próprios erros. A
  linha vermelha marca o acaso. Sob a composição de transformações, nenhum
  modelo distingue acerto de erro.}
  \label{fig:res-auroc}
\end{figure}

A fração na classe mais predita revela a diferença de estilo de falha. Sob
todas as transformações no CIFAR-10, a ResNet-18 colapsa: $76\%$ das predições
em uma única classe no regime de augmentação forte, contra $43\%$ do
Efficient-CapsNet e $38\%$ do DeiT-Tiny (Figura~\ref{fig:res-collapse}). A
Figura~\ref{fig:res-perclass} mostra a forma desse colapso por classe: a
ResNet-18 corrompida acerta $42\%$ de \emph{airplane} e $55\%$ de \emph{bird} e
fica em $0$--$3\%$ em sete das dez classes. O Efficient-CapsNet degrada de
modo mais distribuído, sem por isso acertar mais.

\begin{figure}[t]
  \centering
  \includegraphics[width=\textwidth]{fig08_prediction_collapse_cifar_10.png}
  \caption{CIFAR-10: fração das predições concentradas na classe mais predita.
  Valores próximos de $1/K$ indicam predições ainda distribuídas; valores altos
  indicam colapso.}
  \label{fig:res-collapse}
\end{figure}

\begin{figure}[t]
  \centering
  \includegraphics[width=\textwidth]{fig09_per_class_accuracy_cifar_10.png}
  \caption{CIFAR-10: acurácia por classe e condição de teste, sem augmentação
  de treino, $f = 1{,}00$, agregada sobre sementes.}
  \label{fig:res-perclass}
\end{figure}

Um detalhe da Figura~\ref{fig:res-perclass} merece nota. A classe \emph{frog} é
a única que o Efficient-CapsNet mantém sob rotação grande: $85\%$ limpo e
$85\%$ rotacionado, retenção integral. A ResNet-18 mostra o mesmo padrão nessa
classe ($86\% \to 80\%$). A explicação mais econômica não é arquitetural, e sim
do dado: sapos aparecem em poses e orientações variadas e sem eixo vertical
dominante, ao contrário de cavalos, navios e caminhões. Onde o prior de
orientação é fraco, a rotação não penaliza --- em qualquer arquitetura.

% ---------------------------------------------------------------------
\section{Eficiência de dados}
\label{sec:res:data}

A Tabela~\ref{tab:res-dataeff} traz a estatística $\Delta_f$ da
Equação~\eqref{eq:data-eff}, a queda de acurácia limpa entre $f = 1{,}00$ e
$f = 0{,}33$.

\begin{table}[t]
  \centering
  \caption{Eficiência de dados: $\Delta_f = A_{f=1{,}00} - A_{f=0{,}33}$ em pp
  de acurácia limpa. Valores menores indicam menor dependência do volume de
  dados.}
  \label{tab:res-dataeff}
  \small
  \begin{tabular}{llccc c}
    \toprule
    \multirow{2}{*}{Conjunto} & \multirow{2}{*}{Modelo} & \multicolumn{3}{c}{Augmentação} & \multirow{2}{*}{Média} \\
    \cmidrule(lr){3-5}
     & & nenhuma & padrão & forte & \\
    \midrule
    \multirow{3}{*}{CIFAR-10}
      & Efficient-CapsNet & 10,4 & \phantom{0}6,7 & \phantom{0}6,5 & \textbf{7,9} \\
      & ResNet-18         & 12,6 & \phantom{0}7,5 & \phantom{0}6,5 & 8,9 \\
      & DeiT-Tiny         & \phantom{0}8,1 & 11,0 & 19,0 & 12,7 \\
    \midrule
    \multirow{3}{*}{D-Fire}
      & Efficient-CapsNet & \phantom{0}2,9 & \phantom{0}4,6 & \phantom{0}5,4 & 4,3 \\
      & ResNet-18         & \phantom{0}6,8 & \phantom{0}5,3 & \phantom{0}4,7 & 5,6 \\
      & DeiT-Tiny         & \phantom{0}2,4 & \phantom{0}5,0 & \phantom{0}3,6 & \textbf{3,7} \\
    \bottomrule
  \end{tabular}
\end{table}

O Efficient-CapsNet tem o menor $\Delta_f$ médio no CIFAR-10 ($7{,}9$ contra
$8{,}9$~pp da ResNet-18) e fica $1{,}3$~pp à frente dela no D-Fire. A direção é
a prevista por H2, mas a magnitude não sustenta a hipótese: no teste pareado
por semente e regime de augmentação, a diferença Efficient-CapsNet menos
ResNet-18 é de $-1{,}0$~pp com $p = 0{,}080$ no CIFAR-10 e $-1{,}3$~pp com
$p = 0{,}139$ no D-Fire, ambas com $n = 9$. Nenhuma alcança significância
convencional, e o poder com três sementes é baixo demais para que a ausência de
significância seja informativa em qualquer direção.

Além disso, o ordenamento absoluto não se altera em nenhuma fração: a
ResNet-18 em $f = 0{,}33$ já supera o Efficient-CapsNet em $f = 1{,}00$ no
CIFAR-10 sob augmentação forte ($86{,}8\%$ contra $83{,}0\%$). Uma eventual
inclinação mais suave da curva capsular não se converte em vantagem prática
dentro da faixa investigada.

A Figura~\ref{fig:res-scaling} mostra também um efeito que $\Delta_f$ não
captura: sob rotação grande, mais dados continuam ajudando a ResNet-18
($+4{,}9$~pp de $f = 0{,}33$ a $f = 1{,}00$ com augmentação forte), enquanto
sob todas as transformações as curvas são planas ou levemente decrescentes para
os três modelos. Dados adicionais compram acurácia em distribuição e sob
deslocamento moderado; não compram nada no regime de colapso.

\begin{figure}[t]
  \centering
  \includegraphics[width=\textwidth]{fig05_data_scaling_cifar_10.png}
  \caption{CIFAR-10: acurácia contra fração de dados de treino, por regime de
  augmentação (linhas) e condição de teste (colunas). Barras de erro são o
  desvio padrão entre três sementes.}
  \label{fig:res-scaling}
\end{figure}

% ---------------------------------------------------------------------
\section{Interação com a augmentação}
\label{sec:res:aug}

H3 previa que uma eventual vantagem capsular encolheria ao nível do ruído
quando as linhas de base recebessem augmentação forte. O resultado observado é
mais forte que a hipótese e aponta na direção contrária: não há vantagem
capsular a encolher, e a \emph{desvantagem} capsular \emph{cresce} com a
augmentação.

A Figura~\ref{fig:res-ladder} quantifica o quanto cada arquitetura converte
augmentação em desempenho. Subindo de ``nenhuma'' a ``forte'' no CIFAR-10, a
ResNet-18 ganha $17$~pp em dados limpos e $22$~pp sob rotação; o
Efficient-CapsNet ganha $11$ e $12$~pp; o DeiT-Tiny ganha $1$ e $9$~pp. A
ResNet-18 é, das três, a que melhor aproveita o regime.

\begin{figure}[t]
  \centering
  \includegraphics[width=\textwidth]{fig04_augmentation_ladder_cifar_10.png}
  \caption{CIFAR-10: efeito da intensidade da augmentação de treino sobre cada
  condição de teste, com média sobre as frações de dados. As anotações dão o
  ganho total de ``nenhuma'' a ``forte''.}
  \label{fig:res-ladder}
\end{figure}

A consequência aparece na Figura~\ref{fig:res-heatmap}, que mostra a diferença
entre o Efficient-CapsNet e a melhor linha de base em cada uma das 54
condições do CIFAR-10. Sem augmentação e com $f = 0{,}33$, a diferença é de
$-2{,}3$~pp em dados limpos e chega a $+2{,}3$~pp sob rotação, ambas dentro do
dobro do ruído entre sementes. Com augmentação forte e $f = 1{,}00$, ela é de
$-10{,}3$~pp em dados limpos e $-8{,}7$~pp sob rotação. Em média sobre as
frações, a desvantagem capsular em dados limpos passa de $-2{,}5$~pp (sem
augmentação) para $-8{,}8$ (padrão) e $-13{,}2$~pp (forte).

\begin{figure}[t]
  \centering
  \includegraphics[width=\textwidth]{fig06_advantage_heatmap_cifar_10.png}
  \caption{CIFAR-10: diferença de acurácia entre o Efficient-CapsNet e a melhor
  linha de base em cada condição. Azul indica vantagem capsular, vermelho
  desvantagem; o símbolo $^\circ$ marca células dentro do dobro do ruído entre
  sementes.}
  \label{fig:res-heatmap}
\end{figure}

A interpretação natural é que a invariância que as cápsulas ofereceriam por
construção pode ser adquirida pelas linhas de base por meio de dados, e que
adquiri-la por dados é mais barato e mais eficaz. O resultado é consistente com
a leitura crítica da literatura: o que se credita ao roteamento talvez seja, em
boa parte, atribuível ao objetivo de treino e ao regime de augmentação
(Seção~\ref{sec:protocol:threats}).

Vale registrar o contraponto: nem o Efficient-CapsNet nem o DeiT-Tiny se
beneficiam de augmentação forte em dados limpos; ambos pioram em relação à
augmentação padrão, nos dois conjuntos. Isso
sugere que a perda de margem com reconstrução, ou a geometria das cápsulas
primárias, tolera mal o RandAugment e o Random Erasing --- e que a comparação
sob ``augmentação forte'' talvez não esteja comparando cada arquitetura em seu
melhor. É uma limitação do delineamento, não um resultado sobre cápsulas.

% ---------------------------------------------------------------------
\section{Variabilidade entre sementes}
\label{sec:res:seeds}

A Tabela~\ref{tab:res-seeds} reporta a amplitude entre as três sementes, média
sobre todas as condições de avaliação.

\begin{table}[t]
  \centering
  \caption{Amplitude (máximo menos mínimo) da acurácia entre três sementes, em
  pp, média sobre todas as condições de avaliação.}
  \label{tab:res-seeds}
  \small
  \begin{tabular}{llcccc}
    \toprule
    \multirow{2}{*}{Conjunto} & \multirow{2}{*}{Modelo} & \multicolumn{3}{c}{Augmentação} & \multirow{2}{*}{Média} \\
    \cmidrule(lr){3-5}
     & & nenhuma & padrão & forte & \\
    \midrule
    \multirow{3}{*}{CIFAR-10}
      & Efficient-CapsNet & 0,93 & 0,93 & 0,86 & \textbf{0,91} \\
      & ResNet-18         & 2,01 & 1,17 & 1,06 & 1,42 \\
      & DeiT-Tiny         & 1,72 & 3,39 & 6,75 & 3,95 \\
    \midrule
    \multirow{3}{*}{D-Fire}
      & Efficient-CapsNet & 1,44 & 1,62 & 1,38 & \textbf{1,48} \\
      & ResNet-18         & 1,84 & 1,26 & 1,79 & 1,63 \\
      & DeiT-Tiny         & 2,14 & 7,57 & 7,70 & 5,80 \\
    \bottomrule
  \end{tabular}
\end{table}

Este é o resultado mais nítido a favor do modelo capsular em toda a grade, e
não é sobre robustez nem sobre acurácia. O Efficient-CapsNet é a arquitetura
mais estável entre sementes nos dois conjuntos, com amplitude média de
$0{,}91$~pp no CIFAR-10 --- menos de dois terços da ResNet-18 e menos de um
quarto do DeiT-Tiny. A estabilidade é quase insensível ao regime de
augmentação, enquanto a do DeiT-Tiny se deteriora fortemente com ele (de
$1{,}7$ para $6{,}8$~pp no CIFAR-10).

A amplitude individual mais extrema da grade é de $19{,}4$~pp, no DeiT-Tiny com
augmentação forte e $f = 1{,}00$ no CIFAR-10 --- uma única semente separa
$57\%$ de $76\%$ de acurácia. Os doze piores casos de dispersão da
grade pertencem todos ao DeiT-Tiny, oito deles no D-Fire. Qualquer número reportado para essa arquitetura neste
trabalho deve ser lido com essa dispersão em mente, e a
Figura~\ref{fig:res-seeds} torna explícito que várias comparações envolvendo
o DeiT-Tiny não sobrevivem ao ruído.

\begin{figure}[t]
  \centering
  \includegraphics[width=\textwidth]{fig10_seed_consistency_cifar_10.png}
  \caption{CIFAR-10, $f = 1{,}00$: cada semente desenhada individualmente,
  ligando a acurácia limpa à acurácia sob todas as transformações.}
  \label{fig:res-seeds}
\end{figure}

Uma leitura possível é que a perda de margem com reconstrução, por regularizar
mais fortemente, produz um objetivo com paisagem menos sensível à
inicialização. A hipótese não é testável dentro deste delineamento, que
confunde arquitetura e função objetivo, mas é a explicação candidata a ser
investigada.

% ---------------------------------------------------------------------
\section{Contagem de vitórias}
\label{sec:res:wins}

Somando as 108 condições de avaliação (2 conjuntos $\times$ 3 augmentações
$\times$ 3 frações $\times$ 2 augmentações em teste $\times$ 3 transformações),
a ResNet-18 vence em 57, o DeiT-Tiny em 33 e o Efficient-CapsNet em 18. Das 18
vitórias capsulares, \emph{uma única} excede o dobro do ruído entre sementes:
D-Fire, augmentação forte, $f = 1{,}00$, teste com augmentação forte e todas as
transformações, $59{,}2\%$ contra $57{,}7\%$ da ResNet-18 --- e essa é
justamente uma condição em que ambos os modelos estão a poucos pontos da linha
de base da classe majoritária.

A distribuição das vitórias também localiza os regimes. As 33 vitórias do
DeiT-Tiny concentram-se sob deslocamento (18 delas na composição de todas as
transformações no CIFAR-10, onde todos estão perto do acaso). As 57 vitórias da
ResNet-18 concentram-se em dados limpos e sob rotação, isto é, nos regimes em
que os números significam alguma coisa.

% ---------------------------------------------------------------------
\section{Veredito sobre as hipóteses}
\label{sec:res:verdict}

\begin{table}[t]
  \centering
  \caption{Resumo da avaliação das hipóteses da
  Seção~\ref{sec:protocol:hypotheses}.}
  \label{tab:res-verdict}
  \small
  \begin{tabular}{l p{0.68\textwidth}}
    \toprule
    Hipótese & Veredito \\
    \midrule
    H1 & \textbf{Rejeitada.} Sob rotação grande, a retenção capsular não supera
    a da ResNet-18 no CIFAR-10 ($+1{,}1$~pp, $p = 0{,}13$) e é inferior no
    D-Fire ($-1{,}5$~pp, $p = 0{,}013$). Em acurácia absoluta, a ResNet-18 está
    à frente em seis das nove condições do CIFAR-10, incluindo as seis com
    augmentação. \\
    \addlinespace
    H2 & \textbf{Não sustentada.} A direção prevista aparece ($\Delta_f$ menor
    para as cápsulas em ambos os conjuntos), mas a diferença é de cerca de
    $1$~pp e não atinge significância com três sementes
    ($p = 0{,}080$ e $p = 0{,}139$). O ordenamento absoluto não se inverte em
    nenhuma fração. \\
    \addlinespace
    H3 & \textbf{Rejeitada, com inversão de sinal.} Não há vantagem capsular a
    encolher. A desvantagem em dados limpos no CIFAR-10 \emph{cresce} de
    $-2{,}5$~pp sem augmentação para $-13{,}2$~pp sob augmentação forte. \\
    \bottomrule
  \end{tabular}
\end{table}

O achado central é negativo e deve ser enunciado sem atenuação: nesta grade,
substituir camadas convolucionais por cápsulas com roteamento por
autoatenção não compra robustez a transformações não vistas, nem eficiência de
dados suficiente para compensar a perda de acurácia em distribuição. Onde o
modelo capsular aparece à frente, ou a margem está dentro do ruído entre
sementes, ou os modelos comparados estão no nível do acaso.

Três achados positivos sobrevivem à mesma análise, e nenhum deles é sobre
equivariância:

\begin{description}
\item[Estabilidade.] O Efficient-CapsNet é a arquitetura menos sensível à
semente nos dois conjuntos, com amplitude média de $0{,}91$~pp no CIFAR-10
contra $1{,}42$ da ResNet-18 e $3{,}95$ do DeiT-Tiny
(Seção~\ref{sec:res:seeds}).

\item[Confiança calibrada em distribuição.] A norma do vetor de cápsula é um
escore de detecção de erro competitivo ou superior ao \emph{softmax} em dados
limpos, com AUROC de até $0{,}856$ no CIFAR-10 e $0{,}838$ no D-Fire
(Seção~\ref{sec:res:collapse}).

\item[Competitividade na tarefa binária.] No D-Fire, a diferença para a melhor
linha de base é de $0{,}8$--$3{,}4$~pp conforme o regime, contra
$2{,}5$--$13{,}2$~pp no CIFAR-10. A penalidade capsular é fortemente dependente
da tarefa, o que sugere que conclusões obtidas apenas no CIFAR-10 não
transferem.
\end{description}

% ---------------------------------------------------------------------
\section{Ameaças à validade destes resultados}
\label{sec:res:threats}

Além das ameaças já declaradas na Seção~\ref{sec:protocol:threats}, a análise
deste capítulo acrescenta quatro.

\begin{description}

\item[A composição de corrupções satura.] Sob todas as transformações, os três
modelos ficam no nível do acaso (CIFAR-10) ou da classe majoritária (D-Fire).
Comparações nessa condição não separam robustez; separam estilo de colapso.
Trinta e seis das 108 condições da Seção~\ref{sec:res:wins} pertencem a esse
regime e foram desconsideradas nas conclusões substantivas.

\item[Uma única magnitude de rotação.] A sonda amostra
$\theta \sim \mathcal{U}(-90^{\circ}, 90^{\circ})$ uma vez por imagem. Uma
varredura determinística em $\theta$ produziria uma curva, e é possível que as
cápsulas se saiam melhor em ângulos pequenos e intermediários, onde a
degradação convolucional ainda é suave. O resultado aqui reportado vale para o
regime agregado, não para cada ângulo.

\item[Poder estatístico.] Com três sementes, diferenças abaixo de
aproximadamente $2$~pp são indistinguíveis do ruído na maior parte da grade.
Isso protege contra falsos positivos e impede, simetricamente, afirmar que H2
é falsa: só é possível afirmar que ela não foi demonstrada.

\item[A augmentação forte não é neutra.] Das três arquiteturas, só a ResNet-18
tem a augmentação forte como melhor regime; o Efficient-CapsNet tem a padrão,
nos dois conjuntos. Comparar as três sob augmentação forte é comparar a
ResNet-18 em seu melhor com o Efficient-CapsNet fora dele. A comparação mais favorável ao modelo capsular --- melhor regime
contra melhor regime --- ainda o deixa $8{,}5$~pp atrás no CIFAR-10 e
$1{,}1$~pp atrás no D-Fire, de modo que a conclusão não se inverte; mas as
magnitudes da Seção~\ref{sec:res:aug} são um teto, não uma estimativa pontual.

\end{description}
 a partir do seu .tex principal.
%
%  Este arquivo NÃO carrega pacote nenhum e NÃO declara \documentclass.
%  Ele usa apenas o que o protocolo-experimental.tex já exige:
%    booktabs, multirow, graphicx  (e, se você usar as figuras, um
%    \graphicspath{{figuras/}} no preâmbulo ou o caminho corrigido nos
%    \includegraphics abaixo).
%
%  Se a sua classe for `article`, troque \chapter -> \section e rebaixe
%  os demais níveis.
% =====================================================================

\chapter{Resultados e Discussão}
\label{ch:results}

Este capítulo reporta a grade completa descrita no
Capítulo~\ref{ch:protocol}: 162 treinos (3 arquiteturas $\times$ 2 conjuntos
$\times$ 3 frações de dados $\times$ 3 regimes de augmentação $\times$ 3
sementes), cada um avaliado sob três condições de teste --- limpo, rotação
grande e composição de todas as transformações não vistas --- com e sem
augmentação forte aplicada em teste. São 324 células de resultado e 972
avaliações individuais por semente.

A leitura que segue está organizada em torno de uma decisão metodológica que
precisa ser explicitada antes dos números: \emph{a razão de retenção
$R_c$ da Equação~\eqref{eq:retention} não é usada como métrica principal}.
A Seção~\ref{sec:res:confound} mostra que, nesta grade, $R_c$ é
sistematicamente anticorrelacionada com a acurácia limpa e portanto premia
modelos fracos. As conclusões substantivas repousam sobre a acurácia absoluta
sob deslocamento, lida contra o acaso e contra a linha de base da classe
majoritária.

% ---------------------------------------------------------------------
\section{Desempenho em distribuição}
\label{sec:res:clean}

A Tabela~\ref{tab:res-clean} traz a acurácia de teste limpa. O ordenamento é
estável em toda a grade e não depende da fração de dados.

\begin{table}[t]
  \centering
  \caption{Acurácia de teste limpa (\%) por arquitetura, fração de dados e
  regime de augmentação. Média $\pm$ desvio padrão sobre três sementes.}
  \label{tab:res-clean}
  \small
  \begin{tabular}{llccc}
    \toprule
    \multirow{2}{*}{Modelo} & \multirow{2}{*}{Aug.} & \multicolumn{3}{c}{Fração de dados $f$} \\
    \cmidrule(lr){3-5}
     & & $0{,}33$ & $0{,}66$ & $1{,}00$ \\
    \midrule
    \multicolumn{5}{l}{\textit{CIFAR-10}} \\
    \multirow{3}{*}{Efficient-CapsNet}
      & nenhuma & 63,8 $\pm$ 0,1 & 70,3 $\pm$ 0,4 & 74,2 $\pm$ 0,1 \\
      & padrão  & 78,0 $\pm$ 0,5 & 82,2 $\pm$ 0,1 & 84,8 $\pm$ 0,4 \\
      & forte   & 76,6 $\pm$ 0,3 & 80,7 $\pm$ 0,1 & 83,0 $\pm$ 0,3 \\
    \multirow{3}{*}{ResNet-18}
      & nenhuma & 66,2 $\pm$ 1,9 & 75,3 $\pm$ 0,9 & 78,8 $\pm$ 1,5 \\
      & padrão  & 83,3 $\pm$ 0,5 & 88,6 $\pm$ 0,4 & 90,8 $\pm$ 0,1 \\
      & forte   & \textbf{86,8 $\pm$ 0,6} & \textbf{91,5 $\pm$ 0,1} & \textbf{93,3 $\pm$ 0,1} \\
    \multirow{3}{*}{DeiT-Tiny}
      & nenhuma & 52,8 $\pm$ 0,7 & 58,8 $\pm$ 0,5 & 60,9 $\pm$ 2,2 \\
      & padrão  & 59,2 $\pm$ 1,3 & 61,6 $\pm$ 6,6 & 70,2 $\pm$ 3,5 \\
      & forte   & 49,6 $\pm$ 4,0 & 57,8 $\pm$ 8,0 & 68,7 $\pm$ 10,4 \\
    \midrule
    \multicolumn{5}{l}{\textit{D-Fire}} \\
    \multirow{3}{*}{Efficient-CapsNet}
      & nenhuma & 87,3 $\pm$ 0,5 & 89,2 $\pm$ 0,2 & 90,2 $\pm$ 0,5 \\
      & padrão  & \textbf{89,0 $\pm$ 0,5} & 92,8 $\pm$ 0,4 & 93,6 $\pm$ 0,4 \\
      & forte   & 86,4 $\pm$ 0,0 & 90,3 $\pm$ 0,4 & 91,9 $\pm$ 0,3 \\
    \multirow{3}{*}{ResNet-18}
      & nenhuma & 84,4 $\pm$ 1,7 & 88,7 $\pm$ 0,6 & 91,2 $\pm$ 0,4 \\
      & padrão  & 88,8 $\pm$ 0,9 & 92,7 $\pm$ 0,3 & 94,1 $\pm$ 0,5 \\
      & forte   & \textbf{90,0 $\pm$ 0,7} & \textbf{93,8 $\pm$ 0,2} & \textbf{94,7 $\pm$ 0,2} \\
    \multirow{3}{*}{DeiT-Tiny}
      & nenhuma & 82,2 $\pm$ 0,7 & 84,0 $\pm$ 2,5 & 84,6 $\pm$ 2,9 \\
      & padrão  & 76,0 $\pm$ 7,4 & 78,9 $\pm$ 8,8 & 81,0 $\pm$ 5,4 \\
      & forte   & 73,9 $\pm$ 2,0 & 75,1 $\pm$ 5,3 & 77,5 $\pm$ 8,1 \\
    \bottomrule
  \end{tabular}
\end{table}

No CIFAR-10 a ResNet-18 vence em todas as nove combinações de fração e
augmentação, com margem entre $4{,}4$ e $10{,}5$ pontos percentuais (pp) sobre o
Efficient-CapsNet. O melhor resultado capsular ($84{,}8\%$, augmentação padrão,
$f = 1{,}00$) fica $8{,}5$~pp abaixo do melhor resultado convolucional
($93{,}3\%$). O DeiT-Tiny é o pior dos três por larga margem, o que é esperado:
transformadores de visão treinados do zero, sem pré-treino e com 45.000 imagens
de $32\times32$, operam muito abaixo do seu regime.

No D-Fire o quadro é diferente. A tarefa é binária e mais fácil, e o
Efficient-CapsNet fica a $1{,}1$~pp da ResNet-18 na melhor configuração de cada
um ($93{,}6\%$ contra $94{,}7\%$). Em duas células --- augmentação nenhuma com
$f = 0{,}33$ e $f = 0{,}66$, e augmentação padrão com $f = 0{,}33$ --- o modelo
capsular tem a maior média, mas nenhuma dessas vantagens excede o dobro do
ruído entre sementes.

Duas observações estruturais valem o registro. A primeira: a ResNet-18 é a
única arquitetura que melhora ao subir de augmentação padrão para forte.
Em $f = 1{,}00$, a passagem custa $1{,}8$~pp ao Efficient-CapsNet no CIFAR-10 e
$1{,}7$~pp no D-Fire, e custa $1{,}6$ e $3{,}5$~pp ao DeiT-Tiny, enquanto dá
$2{,}5$ e $0{,}6$~pp à ResNet-18. A segunda: o desvio padrão entre sementes do
Efficient-CapsNet é a menor de todas em quase toda a tabela, ponto retomado na
Seção~\ref{sec:res:seeds}.

% ---------------------------------------------------------------------
\section{Robustez sob transformações não vistas}
\label{sec:res:robust}

A Tabela~\ref{tab:res-robust} reporta a condição $f = 1{,}00$ sob as duas
sondas de deslocamento. Além da acurácia absoluta, listamos $R_c$ e a
\emph{habilidade retida}
\begin{equation}
  S_c = \frac{A^{c} - A^{\text{acaso}}}{A^{\text{clean}} - A^{\text{acaso}}},
  \label{eq:skill}
\end{equation}
que desconta o acerto trivial antes de formar a razão. A diferença entre as
duas colunas é instrutiva: $R_c$ dá crédito por acertar ao acaso, $S_c$ não.

\begin{table}[t]
  \centering
  \caption{Desempenho sob transformações não vistas em $f = 1{,}00$, sem
  augmentação em teste. $A^{c}$ é a acurácia absoluta (\%), $R_c$ a razão de
  retenção (\%) e $S_c$ a habilidade retida (\%) da Equação~\eqref{eq:skill}.
  Acaso: $10\%$ no CIFAR-10, $50\%$ no D-Fire.}
  \label{tab:res-robust}
  \small
  \begin{tabular}{llcccccc}
    \toprule
    \multirow{2}{*}{Modelo} & \multirow{2}{*}{Aug.}
      & \multicolumn{3}{c}{Rotação grande} & \multicolumn{3}{c}{Todas as transformações} \\
    \cmidrule(lr){3-5}\cmidrule(lr){6-8}
     & & $A^{c}$ & $R_c$ & $S_c$ & $A^{c}$ & $R_c$ & $S_c$ \\
    \midrule
    \multicolumn{8}{l}{\textit{CIFAR-10}} \\
    \multirow{3}{*}{Efficient-CapsNet}
      & nenhuma & 38,1 & 51,3 & 43,7 & 12,6 & 17,0 & 4,1 \\
      & padrão  & 52,9 & 62,4 & 57,4 & 12,8 & 15,1 & 3,8 \\
      & forte   & 51,4 & 61,9 & 56,6 & 13,8 & 16,7 & 5,2 \\
    \multirow{3}{*}{ResNet-18}
      & nenhuma & 36,4 & 46,2 & 38,4 & 11,6 & 14,7 & 2,3 \\
      & padrão  & 57,8 & 63,6 & 59,1 & 12,3 & 13,6 & 2,9 \\
      & forte   & \textbf{60,0} & 64,4 & 60,1 & 13,4 & 14,4 & 4,1 \\
    \multirow{3}{*}{DeiT-Tiny}
      & nenhuma & 40,0 & 65,8 & 59,0 & \textbf{16,4} & 26,9 & 12,6 \\
      & padrão  & 52,8 & 75,1 & 71,0 & 15,9 & 22,6 & 9,8 \\
      & forte   & 52,0 & \textbf{75,7} & \textbf{71,6} & \textbf{16,8} & \textbf{24,5} & \textbf{11,6} \\
    \midrule
    \multicolumn{8}{l}{\textit{D-Fire}} \\
    \multirow{3}{*}{Efficient-CapsNet}
      & nenhuma & 61,4 & 68,1 & 28,4 & 53,1 & 58,9 & 7,7 \\
      & padrão  & 77,6 & 82,9 & 63,2 & 56,1 & 59,9 & 13,9 \\
      & forte   & 80,0 & 87,0 & 71,5 & 60,5 & 65,8 & 25,0 \\
    \multirow{3}{*}{ResNet-18}
      & nenhuma & 63,4 & 69,5 & 32,4 & 52,5 & 57,5 & 6,0 \\
      & padrão  & 78,7 & 83,7 & 65,2 & 54,9 & 58,4 & 11,2 \\
      & forte   & \textbf{80,9} & 85,4 & 69,1 & \textbf{60,7} & 64,1 & 23,9 \\
    \multirow{3}{*}{DeiT-Tiny}
      & nenhuma & 63,5 & 75,0 & 38,9 & 52,9 & 62,5 & 8,4 \\
      & padrão  & 78,4 & 96,9 & 91,8 & 53,7 & 66,4 & 12,1 \\
      & forte   & 76,0 & \textbf{98,0} & \textbf{94,4} & 59,9 & \textbf{77,2} & \textbf{35,9} \\
    \bottomrule
  \end{tabular}
\end{table}

\subsection{A sonda de rotação}

A rotação grande é o teste direto de H1, e o resultado é desfavorável à
hipótese. No CIFAR-10 com augmentação padrão ou forte, a ResNet-18 tem a maior
acurácia absoluta sob rotação --- $57{,}8\%$ e $60{,}0\%$ contra $52{,}9\%$ e
$51{,}4\%$ do Efficient-CapsNet. A vantagem convolucional sob rotação
($+8{,}6$~pp em $f = 1{,}00$, augmentação forte) é maior, em valor absoluto, do
que qualquer vantagem capsular observada em toda a grade.

Apenas no regime sem augmentação nenhuma o Efficient-CapsNet fica à frente sob
rotação, e por pouco: $38{,}1\%$ contra $36{,}4\%$ da ResNet-18. Em teste
pareado por semente sobre todas as 27 células do CIFAR-10, a diferença de
retenção entre Efficient-CapsNet e ResNet-18 sob rotação é de $+1{,}1$~pp com
$p = 0{,}13$ --- indistinguível de zero. No D-Fire a mesma comparação dá
$-1{,}5$~pp com $p = 0{,}013$, ou seja, uma desvantagem capsular
estatisticamente detectável. Contra o DeiT-Tiny a diferença de retenção é de
$-15{,}1$~pp no CIFAR-10 e $-11{,}7$~pp no D-Fire, ambas com $p < 10^{-12}$.

A Figura~\ref{fig:res-dumbbell} torna o padrão visível: as barras do
Efficient-CapsNet e da ResNet-18 partem de pontos diferentes e caem
aproximadamente na mesma proporção, sem que a distância entre elas se feche.

\begin{figure}[t]
  \centering
  \includegraphics[width=\textwidth]{fig01_shift_dumbbell_cifar_10.png}
  \caption{CIFAR-10: acurácia limpa (círculo aberto) e acurácia sob
  deslocamento (marcador preenchido) para cada combinação de arquitetura,
  fração de dados e regime de augmentação. O comprimento da barra é a
  degradação absoluta.}
  \label{fig:res-dumbbell}
\end{figure}

\subsection{A composição de todas as transformações é um piso, não uma medida}
\label{sec:res:floor}

Sob a composição das seis corrupções, os três modelos deixam de funcionar. No
CIFAR-10 a acurácia cai para a faixa de $11{,}1\%$ a $19{,}1\%$, contra
$10{,}0\%$ de acaso, em todas as 54 células. No D-Fire os valores ficam entre
$-2{,}2$ e $+7{,}8$~pp da linha de base da classe majoritária ($53{,}4\%$); sem
augmentação de treino, 14 das 18 células ficam \emph{abaixo} dela, isto é,
perde para um classificador que ignora a imagem
(Figura~\ref{fig:res-baselines}).

Isso muda o que se pode concluir a partir dessa condição. Comparar retenção
entre modelos que estão todos no nível do acaso compara padrões de colapso, não
robustez. As diferenças de $R_c$ que aparecem ali --- $26{,}9\%$ para o
DeiT-Tiny contra $14{,}7\%$ para a ResNet-18, por exemplo --- correspondem, em
acurácia absoluta, a $16{,}4\%$ contra $11{,}6\%$: quatro pontos e meio acima e
um ponto e meio acima do acaso, respectivamente. Nenhum dos dois é utilizável.

\begin{figure}[t]
  \centering
  \includegraphics[width=\textwidth]{fig11_cross_dataset_vs_baseline.png}
  \caption{Os dois conjuntos contra suas próprias linhas de base triviais.
  Barras abaixo da linha pontilhada perdem para um classificador que ignora a
  imagem. Sob todas as transformações, nenhuma arquitetura se sustenta em
  nenhum dos conjuntos.}
  \label{fig:res-baselines}
\end{figure}

A recomendação que decorre daí é metodológica e vale além deste trabalho:
composições de corrupções só informam enquanto os modelos comparados
permanecerem acima do acaso por margem apreciável. Abaixo disso, a métrica
mede outra coisa. A Seção~\ref{sec:res:collapse} identifica exatamente o que.

\subsection{Por que a retenção não pode ser a métrica principal}
\label{sec:res:confound}

O DeiT-Tiny tem a maior razão de retenção em praticamente toda a
Tabela~\ref{tab:res-robust} e a menor acurácia limpa em toda a
Tabela~\ref{tab:res-clean}. Não é coincidência. A Figura~\ref{fig:res-confound}
mostra a correlação entre acurácia limpa e retenção no CIFAR-10:
$r = -0{,}44$ sob rotação e $r = -0{,}91$ sob todas as transformações. As
inclinações ajustadas são de $-34$ e $-48$~pp de retenção por unidade de
acurácia limpa.

\begin{figure}[t]
  \centering
  \includegraphics[width=\textwidth]{fig03_retention_confound_cifar_10.png}
  \caption{CIFAR-10: razão de retenção contra acurácia limpa. A relação
  negativa significa que a retenção, isolada, recompensa modelos que já eram
  ruins em distribuição.}
  \label{fig:res-confound}
\end{figure}

O mecanismo é aritmético. A acurácia sob deslocamento tem piso no acaso,
enquanto a acurácia limpa não tem teto próximo; um modelo que sai de $60\%$ e
cai para $40\%$ retém $67\%$, e um que sai de $93\%$ e cai para $60\%$ retém
$64\%$ --- embora o segundo esteja $20$~pp à frente onde importa. A
Figura~\ref{fig:res-frontier} mostra a mesma informação no plano
(acurácia limpa, acurácia deslocada): sob rotação a correlação entre os dois
eixos é $+0{,}68$, e sob todas as transformações ela \emph{inverte} de sinal
($-0{,}74$), que é a assinatura do regime de colapso descrito acima.

\begin{figure}[t]
  \centering
  \includegraphics[width=\textwidth]{fig02_robustness_frontier_cifar_10.png}
  \caption{CIFAR-10: acurácia sob deslocamento contra acurácia limpa. A
  diagonal tracejada marca invariância perfeita. Separação vertical entre cores
  na mesma abscissa seria o efeito arquitetural puro; ela é pequena sob rotação
  e some sob todas as transformações.}
  \label{fig:res-frontier}
\end{figure}

Por isso a Tabela~\ref{tab:res-robust} reporta $A^c$ em primeiro lugar, e por
isso H1 é avaliada sobre acurácia absoluta e sobre retenção pareada por
semente, não sobre médias de $R_c$.

% ---------------------------------------------------------------------
\section{Colapso de predição e detecção de erro}
\label{sec:res:collapse}

Dois diagnósticos separam ``degradação graciosa'' de ``falha silenciosa'': a
fração de predições que caem na classe mais predita (\emph{top-pred share}) e a
AUROC da confiança do modelo como detector dos próprios erros. A
Tabela~\ref{tab:res-auroc} traz os dois.

\begin{table}[t]
  \centering
  \caption{AUROC de detecção de erro (a confiança separa acertos de erros?) e
  fração de predições na classe mais predita, médias sobre as três frações de
  dados, sem augmentação em teste. AUROC de $0{,}50$ é o acaso; \emph{share} de
  $0{,}10$ (CIFAR-10) e $0{,}53$ (D-Fire) corresponde a predições
  balanceadas.}
  \label{tab:res-auroc}
  \footnotesize
  \setlength{\tabcolsep}{4pt}
  \begin{tabular}{lccccccccc}
    \toprule
    \multirow{2}{*}{Modelo}
      & \multicolumn{3}{c}{Limpo} & \multicolumn{3}{c}{Rotação grande} & \multicolumn{3}{c}{Todas} \\
    \cmidrule(lr){2-4}\cmidrule(lr){5-7}\cmidrule(lr){8-10}
     & nenh. & padr. & forte & nenh. & padr. & forte & nenh. & padr. & forte \\
    \midrule
    \multicolumn{10}{l}{\textit{CIFAR-10 --- AUROC}} \\
    Efficient-CapsNet & 0,799 & 0,856 & 0,848 & 0,707 & 0,762 & 0,766 & 0,533 & 0,547 & 0,548 \\
    ResNet-18         & 0,788 & 0,851 & 0,867 & 0,697 & 0,790 & 0,804 & 0,540 & 0,527 & 0,540 \\
    DeiT-Tiny         & 0,716 & 0,736 & 0,720 & 0,677 & 0,722 & 0,713 & 0,560 & 0,587 & 0,618 \\
    \addlinespace
    \multicolumn{10}{l}{\textit{CIFAR-10 --- fração na classe mais predita}} \\
    Efficient-CapsNet & 0,12 & 0,11 & 0,11 & 0,24 & 0,18 & 0,20 & 0,39 & 0,48 & 0,43 \\
    ResNet-18         & 0,13 & 0,11 & 0,11 & 0,28 & 0,18 & 0,18 & 0,55 & 0,70 & 0,76 \\
    DeiT-Tiny         & 0,13 & 0,12 & 0,13 & 0,21 & 0,15 & 0,16 & 0,38 & 0,48 & 0,38 \\
    \midrule
    \multicolumn{10}{l}{\textit{D-Fire --- AUROC}} \\
    Efficient-CapsNet & 0,837 & 0,838 & 0,836 & 0,581 & 0,696 & 0,728 & 0,515 & 0,558 & 0,595 \\
    ResNet-18         & 0,826 & 0,845 & 0,881 & 0,610 & 0,736 & 0,745 & 0,532 & 0,562 & 0,601 \\
    DeiT-Tiny         & 0,769 & 0,755 & 0,731 & 0,627 & 0,727 & 0,726 & 0,525 & 0,560 & 0,566 \\
    \addlinespace
    \multicolumn{10}{l}{\textit{D-Fire --- fração na classe mais predita}} \\
    Efficient-CapsNet & 0,56 & 0,54 & 0,55 & 0,78 & 0,64 & 0,62 & 0,75 & 0,60 & 0,58 \\
    ResNet-18         & 0,53 & 0,53 & 0,53 & 0,70 & 0,61 & 0,62 & 0,65 & 0,58 & 0,62 \\
    DeiT-Tiny         & 0,55 & 0,54 & 0,54 & 0,55 & 0,54 & 0,54 & 0,70 & 0,63 & 0,63 \\
    \bottomrule
  \end{tabular}
\end{table}

Em dados limpos, o Efficient-CapsNet tem a melhor detecção de erro do CIFAR-10
sob augmentação padrão ($0{,}856$ contra $0{,}851$ da ResNet-18) e a melhor do
D-Fire no regime sem augmentação ($0{,}837$ contra $0{,}826$), ficando a menos
de um centésimo da ResNet-18 sob augmentação padrão. A norma do vetor de cápsula, portanto,
funciona como escore de confiança pelo menos tão bem quanto o \emph{softmax},
o que é um resultado positivo para o modelo capsular e independente da
acurácia.

Sob deslocamento essa propriedade não sobrevive. Sob rotação a AUROC capsular
cai para $0{,}707$--$0{,}766$ no CIFAR-10 e para $0{,}581$ no D-Fire sem
augmentação --- neste último caso, com $77{,}6\%$ das predições em uma única
classe. Ou seja: o modelo rotacionado no D-Fire responde quase sempre a mesma
coisa e não sabe que está errando. Sob todas as transformações, a AUROC dos
três modelos fica entre $0{,}515$ e $0{,}618$, praticamente o acaso
(Figura~\ref{fig:res-auroc}). A falha, nessa condição, é silenciosa em todas as
arquiteturas.

\begin{figure}[t]
  \centering
  \includegraphics[width=\textwidth]{fig07_error_detection_auroc_cifar_10.png}
  \caption{CIFAR-10: AUROC da confiança como detector dos próprios erros. A
  linha vermelha marca o acaso. Sob a composição de transformações, nenhum
  modelo distingue acerto de erro.}
  \label{fig:res-auroc}
\end{figure}

A fração na classe mais predita revela a diferença de estilo de falha. Sob
todas as transformações no CIFAR-10, a ResNet-18 colapsa: $76\%$ das predições
em uma única classe no regime de augmentação forte, contra $43\%$ do
Efficient-CapsNet e $38\%$ do DeiT-Tiny (Figura~\ref{fig:res-collapse}). A
Figura~\ref{fig:res-perclass} mostra a forma desse colapso por classe: a
ResNet-18 corrompida acerta $42\%$ de \emph{airplane} e $55\%$ de \emph{bird} e
fica em $0$--$3\%$ em sete das dez classes. O Efficient-CapsNet degrada de
modo mais distribuído, sem por isso acertar mais.

\begin{figure}[t]
  \centering
  \includegraphics[width=\textwidth]{fig08_prediction_collapse_cifar_10.png}
  \caption{CIFAR-10: fração das predições concentradas na classe mais predita.
  Valores próximos de $1/K$ indicam predições ainda distribuídas; valores altos
  indicam colapso.}
  \label{fig:res-collapse}
\end{figure}

\begin{figure}[t]
  \centering
  \includegraphics[width=\textwidth]{fig09_per_class_accuracy_cifar_10.png}
  \caption{CIFAR-10: acurácia por classe e condição de teste, sem augmentação
  de treino, $f = 1{,}00$, agregada sobre sementes.}
  \label{fig:res-perclass}
\end{figure}

Um detalhe da Figura~\ref{fig:res-perclass} merece nota. A classe \emph{frog} é
a única que o Efficient-CapsNet mantém sob rotação grande: $85\%$ limpo e
$85\%$ rotacionado, retenção integral. A ResNet-18 mostra o mesmo padrão nessa
classe ($86\% \to 80\%$). A explicação mais econômica não é arquitetural, e sim
do dado: sapos aparecem em poses e orientações variadas e sem eixo vertical
dominante, ao contrário de cavalos, navios e caminhões. Onde o prior de
orientação é fraco, a rotação não penaliza --- em qualquer arquitetura.

% ---------------------------------------------------------------------
\section{Eficiência de dados}
\label{sec:res:data}

A Tabela~\ref{tab:res-dataeff} traz a estatística $\Delta_f$ da
Equação~\eqref{eq:data-eff}, a queda de acurácia limpa entre $f = 1{,}00$ e
$f = 0{,}33$.

\begin{table}[t]
  \centering
  \caption{Eficiência de dados: $\Delta_f = A_{f=1{,}00} - A_{f=0{,}33}$ em pp
  de acurácia limpa. Valores menores indicam menor dependência do volume de
  dados.}
  \label{tab:res-dataeff}
  \small
  \begin{tabular}{llccc c}
    \toprule
    \multirow{2}{*}{Conjunto} & \multirow{2}{*}{Modelo} & \multicolumn{3}{c}{Augmentação} & \multirow{2}{*}{Média} \\
    \cmidrule(lr){3-5}
     & & nenhuma & padrão & forte & \\
    \midrule
    \multirow{3}{*}{CIFAR-10}
      & Efficient-CapsNet & 10,4 & \phantom{0}6,7 & \phantom{0}6,5 & \textbf{7,9} \\
      & ResNet-18         & 12,6 & \phantom{0}7,5 & \phantom{0}6,5 & 8,9 \\
      & DeiT-Tiny         & \phantom{0}8,1 & 11,0 & 19,0 & 12,7 \\
    \midrule
    \multirow{3}{*}{D-Fire}
      & Efficient-CapsNet & \phantom{0}2,9 & \phantom{0}4,6 & \phantom{0}5,4 & 4,3 \\
      & ResNet-18         & \phantom{0}6,8 & \phantom{0}5,3 & \phantom{0}4,7 & 5,6 \\
      & DeiT-Tiny         & \phantom{0}2,4 & \phantom{0}5,0 & \phantom{0}3,6 & \textbf{3,7} \\
    \bottomrule
  \end{tabular}
\end{table}

O Efficient-CapsNet tem o menor $\Delta_f$ médio no CIFAR-10 ($7{,}9$ contra
$8{,}9$~pp da ResNet-18) e fica $1{,}3$~pp à frente dela no D-Fire. A direção é
a prevista por H2, mas a magnitude não sustenta a hipótese: no teste pareado
por semente e regime de augmentação, a diferença Efficient-CapsNet menos
ResNet-18 é de $-1{,}0$~pp com $p = 0{,}080$ no CIFAR-10 e $-1{,}3$~pp com
$p = 0{,}139$ no D-Fire, ambas com $n = 9$. Nenhuma alcança significância
convencional, e o poder com três sementes é baixo demais para que a ausência de
significância seja informativa em qualquer direção.

Além disso, o ordenamento absoluto não se altera em nenhuma fração: a
ResNet-18 em $f = 0{,}33$ já supera o Efficient-CapsNet em $f = 1{,}00$ no
CIFAR-10 sob augmentação forte ($86{,}8\%$ contra $83{,}0\%$). Uma eventual
inclinação mais suave da curva capsular não se converte em vantagem prática
dentro da faixa investigada.

A Figura~\ref{fig:res-scaling} mostra também um efeito que $\Delta_f$ não
captura: sob rotação grande, mais dados continuam ajudando a ResNet-18
($+4{,}9$~pp de $f = 0{,}33$ a $f = 1{,}00$ com augmentação forte), enquanto
sob todas as transformações as curvas são planas ou levemente decrescentes para
os três modelos. Dados adicionais compram acurácia em distribuição e sob
deslocamento moderado; não compram nada no regime de colapso.

\begin{figure}[t]
  \centering
  \includegraphics[width=\textwidth]{fig05_data_scaling_cifar_10.png}
  \caption{CIFAR-10: acurácia contra fração de dados de treino, por regime de
  augmentação (linhas) e condição de teste (colunas). Barras de erro são o
  desvio padrão entre três sementes.}
  \label{fig:res-scaling}
\end{figure}

% ---------------------------------------------------------------------
\section{Interação com a augmentação}
\label{sec:res:aug}

H3 previa que uma eventual vantagem capsular encolheria ao nível do ruído
quando as linhas de base recebessem augmentação forte. O resultado observado é
mais forte que a hipótese e aponta na direção contrária: não há vantagem
capsular a encolher, e a \emph{desvantagem} capsular \emph{cresce} com a
augmentação.

A Figura~\ref{fig:res-ladder} quantifica o quanto cada arquitetura converte
augmentação em desempenho. Subindo de ``nenhuma'' a ``forte'' no CIFAR-10, a
ResNet-18 ganha $17$~pp em dados limpos e $22$~pp sob rotação; o
Efficient-CapsNet ganha $11$ e $12$~pp; o DeiT-Tiny ganha $1$ e $9$~pp. A
ResNet-18 é, das três, a que melhor aproveita o regime.

\begin{figure}[t]
  \centering
  \includegraphics[width=\textwidth]{fig04_augmentation_ladder_cifar_10.png}
  \caption{CIFAR-10: efeito da intensidade da augmentação de treino sobre cada
  condição de teste, com média sobre as frações de dados. As anotações dão o
  ganho total de ``nenhuma'' a ``forte''.}
  \label{fig:res-ladder}
\end{figure}

A consequência aparece na Figura~\ref{fig:res-heatmap}, que mostra a diferença
entre o Efficient-CapsNet e a melhor linha de base em cada uma das 54
condições do CIFAR-10. Sem augmentação e com $f = 0{,}33$, a diferença é de
$-2{,}3$~pp em dados limpos e chega a $+2{,}3$~pp sob rotação, ambas dentro do
dobro do ruído entre sementes. Com augmentação forte e $f = 1{,}00$, ela é de
$-10{,}3$~pp em dados limpos e $-8{,}7$~pp sob rotação. Em média sobre as
frações, a desvantagem capsular em dados limpos passa de $-2{,}5$~pp (sem
augmentação) para $-8{,}8$ (padrão) e $-13{,}2$~pp (forte).

\begin{figure}[t]
  \centering
  \includegraphics[width=\textwidth]{fig06_advantage_heatmap_cifar_10.png}
  \caption{CIFAR-10: diferença de acurácia entre o Efficient-CapsNet e a melhor
  linha de base em cada condição. Azul indica vantagem capsular, vermelho
  desvantagem; o símbolo $^\circ$ marca células dentro do dobro do ruído entre
  sementes.}
  \label{fig:res-heatmap}
\end{figure}

A interpretação natural é que a invariância que as cápsulas ofereceriam por
construção pode ser adquirida pelas linhas de base por meio de dados, e que
adquiri-la por dados é mais barato e mais eficaz. O resultado é consistente com
a leitura crítica da literatura: o que se credita ao roteamento talvez seja, em
boa parte, atribuível ao objetivo de treino e ao regime de augmentação
(Seção~\ref{sec:protocol:threats}).

Vale registrar o contraponto: nem o Efficient-CapsNet nem o DeiT-Tiny se
beneficiam de augmentação forte em dados limpos; ambos pioram em relação à
augmentação padrão, nos dois conjuntos. Isso
sugere que a perda de margem com reconstrução, ou a geometria das cápsulas
primárias, tolera mal o RandAugment e o Random Erasing --- e que a comparação
sob ``augmentação forte'' talvez não esteja comparando cada arquitetura em seu
melhor. É uma limitação do delineamento, não um resultado sobre cápsulas.

% ---------------------------------------------------------------------
\section{Variabilidade entre sementes}
\label{sec:res:seeds}

A Tabela~\ref{tab:res-seeds} reporta a amplitude entre as três sementes, média
sobre todas as condições de avaliação.

\begin{table}[t]
  \centering
  \caption{Amplitude (máximo menos mínimo) da acurácia entre três sementes, em
  pp, média sobre todas as condições de avaliação.}
  \label{tab:res-seeds}
  \small
  \begin{tabular}{llcccc}
    \toprule
    \multirow{2}{*}{Conjunto} & \multirow{2}{*}{Modelo} & \multicolumn{3}{c}{Augmentação} & \multirow{2}{*}{Média} \\
    \cmidrule(lr){3-5}
     & & nenhuma & padrão & forte & \\
    \midrule
    \multirow{3}{*}{CIFAR-10}
      & Efficient-CapsNet & 0,93 & 0,93 & 0,86 & \textbf{0,91} \\
      & ResNet-18         & 2,01 & 1,17 & 1,06 & 1,42 \\
      & DeiT-Tiny         & 1,72 & 3,39 & 6,75 & 3,95 \\
    \midrule
    \multirow{3}{*}{D-Fire}
      & Efficient-CapsNet & 1,44 & 1,62 & 1,38 & \textbf{1,48} \\
      & ResNet-18         & 1,84 & 1,26 & 1,79 & 1,63 \\
      & DeiT-Tiny         & 2,14 & 7,57 & 7,70 & 5,80 \\
    \bottomrule
  \end{tabular}
\end{table}

Este é o resultado mais nítido a favor do modelo capsular em toda a grade, e
não é sobre robustez nem sobre acurácia. O Efficient-CapsNet é a arquitetura
mais estável entre sementes nos dois conjuntos, com amplitude média de
$0{,}91$~pp no CIFAR-10 --- menos de dois terços da ResNet-18 e menos de um
quarto do DeiT-Tiny. A estabilidade é quase insensível ao regime de
augmentação, enquanto a do DeiT-Tiny se deteriora fortemente com ele (de
$1{,}7$ para $6{,}8$~pp no CIFAR-10).

A amplitude individual mais extrema da grade é de $19{,}4$~pp, no DeiT-Tiny com
augmentação forte e $f = 1{,}00$ no CIFAR-10 --- uma única semente separa
$57\%$ de $76\%$ de acurácia. Os doze piores casos de dispersão da
grade pertencem todos ao DeiT-Tiny, oito deles no D-Fire. Qualquer número reportado para essa arquitetura neste
trabalho deve ser lido com essa dispersão em mente, e a
Figura~\ref{fig:res-seeds} torna explícito que várias comparações envolvendo
o DeiT-Tiny não sobrevivem ao ruído.

\begin{figure}[t]
  \centering
  \includegraphics[width=\textwidth]{fig10_seed_consistency_cifar_10.png}
  \caption{CIFAR-10, $f = 1{,}00$: cada semente desenhada individualmente,
  ligando a acurácia limpa à acurácia sob todas as transformações.}
  \label{fig:res-seeds}
\end{figure}

Uma leitura possível é que a perda de margem com reconstrução, por regularizar
mais fortemente, produz um objetivo com paisagem menos sensível à
inicialização. A hipótese não é testável dentro deste delineamento, que
confunde arquitetura e função objetivo, mas é a explicação candidata a ser
investigada.

% ---------------------------------------------------------------------
\section{Contagem de vitórias}
\label{sec:res:wins}

Somando as 108 condições de avaliação (2 conjuntos $\times$ 3 augmentações
$\times$ 3 frações $\times$ 2 augmentações em teste $\times$ 3 transformações),
a ResNet-18 vence em 57, o DeiT-Tiny em 33 e o Efficient-CapsNet em 18. Das 18
vitórias capsulares, \emph{uma única} excede o dobro do ruído entre sementes:
D-Fire, augmentação forte, $f = 1{,}00$, teste com augmentação forte e todas as
transformações, $59{,}2\%$ contra $57{,}7\%$ da ResNet-18 --- e essa é
justamente uma condição em que ambos os modelos estão a poucos pontos da linha
de base da classe majoritária.

A distribuição das vitórias também localiza os regimes. As 33 vitórias do
DeiT-Tiny concentram-se sob deslocamento (18 delas na composição de todas as
transformações no CIFAR-10, onde todos estão perto do acaso). As 57 vitórias da
ResNet-18 concentram-se em dados limpos e sob rotação, isto é, nos regimes em
que os números significam alguma coisa.

% ---------------------------------------------------------------------
\section{Veredito sobre as hipóteses}
\label{sec:res:verdict}

\begin{table}[t]
  \centering
  \caption{Resumo da avaliação das hipóteses da
  Seção~\ref{sec:protocol:hypotheses}.}
  \label{tab:res-verdict}
  \small
  \begin{tabular}{l p{0.68\textwidth}}
    \toprule
    Hipótese & Veredito \\
    \midrule
    H1 & \textbf{Rejeitada.} Sob rotação grande, a retenção capsular não supera
    a da ResNet-18 no CIFAR-10 ($+1{,}1$~pp, $p = 0{,}13$) e é inferior no
    D-Fire ($-1{,}5$~pp, $p = 0{,}013$). Em acurácia absoluta, a ResNet-18 está
    à frente em seis das nove condições do CIFAR-10, incluindo as seis com
    augmentação. \\
    \addlinespace
    H2 & \textbf{Não sustentada.} A direção prevista aparece ($\Delta_f$ menor
    para as cápsulas em ambos os conjuntos), mas a diferença é de cerca de
    $1$~pp e não atinge significância com três sementes
    ($p = 0{,}080$ e $p = 0{,}139$). O ordenamento absoluto não se inverte em
    nenhuma fração. \\
    \addlinespace
    H3 & \textbf{Rejeitada, com inversão de sinal.} Não há vantagem capsular a
    encolher. A desvantagem em dados limpos no CIFAR-10 \emph{cresce} de
    $-2{,}5$~pp sem augmentação para $-13{,}2$~pp sob augmentação forte. \\
    \bottomrule
  \end{tabular}
\end{table}

O achado central é negativo e deve ser enunciado sem atenuação: nesta grade,
substituir camadas convolucionais por cápsulas com roteamento por
autoatenção não compra robustez a transformações não vistas, nem eficiência de
dados suficiente para compensar a perda de acurácia em distribuição. Onde o
modelo capsular aparece à frente, ou a margem está dentro do ruído entre
sementes, ou os modelos comparados estão no nível do acaso.

Três achados positivos sobrevivem à mesma análise, e nenhum deles é sobre
equivariância:

\begin{description}
\item[Estabilidade.] O Efficient-CapsNet é a arquitetura menos sensível à
semente nos dois conjuntos, com amplitude média de $0{,}91$~pp no CIFAR-10
contra $1{,}42$ da ResNet-18 e $3{,}95$ do DeiT-Tiny
(Seção~\ref{sec:res:seeds}).

\item[Confiança calibrada em distribuição.] A norma do vetor de cápsula é um
escore de detecção de erro competitivo ou superior ao \emph{softmax} em dados
limpos, com AUROC de até $0{,}856$ no CIFAR-10 e $0{,}838$ no D-Fire
(Seção~\ref{sec:res:collapse}).

\item[Competitividade na tarefa binária.] No D-Fire, a diferença para a melhor
linha de base é de $0{,}8$--$3{,}4$~pp conforme o regime, contra
$2{,}5$--$13{,}2$~pp no CIFAR-10. A penalidade capsular é fortemente dependente
da tarefa, o que sugere que conclusões obtidas apenas no CIFAR-10 não
transferem.
\end{description}

% ---------------------------------------------------------------------
\section{Ameaças à validade destes resultados}
\label{sec:res:threats}

Além das ameaças já declaradas na Seção~\ref{sec:protocol:threats}, a análise
deste capítulo acrescenta quatro.

\begin{description}

\item[A composição de corrupções satura.] Sob todas as transformações, os três
modelos ficam no nível do acaso (CIFAR-10) ou da classe majoritária (D-Fire).
Comparações nessa condição não separam robustez; separam estilo de colapso.
Trinta e seis das 108 condições da Seção~\ref{sec:res:wins} pertencem a esse
regime e foram desconsideradas nas conclusões substantivas.

\item[Uma única magnitude de rotação.] A sonda amostra
$\theta \sim \mathcal{U}(-90^{\circ}, 90^{\circ})$ uma vez por imagem. Uma
varredura determinística em $\theta$ produziria uma curva, e é possível que as
cápsulas se saiam melhor em ângulos pequenos e intermediários, onde a
degradação convolucional ainda é suave. O resultado aqui reportado vale para o
regime agregado, não para cada ângulo.

\item[Poder estatístico.] Com três sementes, diferenças abaixo de
aproximadamente $2$~pp são indistinguíveis do ruído na maior parte da grade.
Isso protege contra falsos positivos e impede, simetricamente, afirmar que H2
é falsa: só é possível afirmar que ela não foi demonstrada.

\item[A augmentação forte não é neutra.] Das três arquiteturas, só a ResNet-18
tem a augmentação forte como melhor regime; o Efficient-CapsNet tem a padrão,
nos dois conjuntos. Comparar as três sob augmentação forte é comparar a
ResNet-18 em seu melhor com o Efficient-CapsNet fora dele. A comparação mais favorável ao modelo capsular --- melhor regime
contra melhor regime --- ainda o deixa $8{,}5$~pp atrás no CIFAR-10 e
$1{,}1$~pp atrás no D-Fire, de modo que a conclusão não se inverte; mas as
magnitudes da Seção~\ref{sec:res:aug} são um teto, não uma estimativa pontual.

\end{description}
 a partir do seu .tex principal.
%
%  Este arquivo NÃO carrega pacote nenhum e NÃO declara \documentclass.
%  Ele usa apenas o que o protocolo-experimental.tex já exige:
%    booktabs, multirow, graphicx  (e, se você usar as figuras, um
%    \graphicspath{{figuras/}} no preâmbulo ou o caminho corrigido nos
%    \includegraphics abaixo).
%
%  Se a sua classe for `article`, troque \chapter -> \section e rebaixe
%  os demais níveis.
% =====================================================================

\chapter{Resultados e Discussão}
\label{ch:results}

Este capítulo reporta a grade completa descrita no
Capítulo~\ref{ch:protocol}: 162 treinos (3 arquiteturas $\times$ 2 conjuntos
$\times$ 3 frações de dados $\times$ 3 regimes de augmentação $\times$ 3
sementes), cada um avaliado sob três condições de teste --- limpo, rotação
grande e composição de todas as transformações não vistas --- com e sem
augmentação forte aplicada em teste. São 324 células de resultado e 972
avaliações individuais por semente.

A leitura que segue está organizada em torno de uma decisão metodológica que
precisa ser explicitada antes dos números: \emph{a razão de retenção
$R_c$ da Equação~\eqref{eq:retention} não é usada como métrica principal}.
A Seção~\ref{sec:res:confound} mostra que, nesta grade, $R_c$ é
sistematicamente anticorrelacionada com a acurácia limpa e portanto premia
modelos fracos. As conclusões substantivas repousam sobre a acurácia absoluta
sob deslocamento, lida contra o acaso e contra a linha de base da classe
majoritária.

% ---------------------------------------------------------------------
\section{Desempenho em distribuição}
\label{sec:res:clean}

A Tabela~\ref{tab:res-clean} traz a acurácia de teste limpa. O ordenamento é
estável em toda a grade e não depende da fração de dados.

\begin{table}[t]
  \centering
  \caption{Acurácia de teste limpa (\%) por arquitetura, fração de dados e
  regime de augmentação. Média $\pm$ desvio padrão sobre três sementes.}
  \label{tab:res-clean}
  \small
  \begin{tabular}{llccc}
    \toprule
    \multirow{2}{*}{Modelo} & \multirow{2}{*}{Aug.} & \multicolumn{3}{c}{Fração de dados $f$} \\
    \cmidrule(lr){3-5}
     & & $0{,}33$ & $0{,}66$ & $1{,}00$ \\
    \midrule
    \multicolumn{5}{l}{\textit{CIFAR-10}} \\
    \multirow{3}{*}{Efficient-CapsNet}
      & nenhuma & 63,8 $\pm$ 0,1 & 70,3 $\pm$ 0,4 & 74,2 $\pm$ 0,1 \\
      & padrão  & 78,0 $\pm$ 0,5 & 82,2 $\pm$ 0,1 & 84,8 $\pm$ 0,4 \\
      & forte   & 76,6 $\pm$ 0,3 & 80,7 $\pm$ 0,1 & 83,0 $\pm$ 0,3 \\
    \multirow{3}{*}{ResNet-18}
      & nenhuma & 66,2 $\pm$ 1,9 & 75,3 $\pm$ 0,9 & 78,8 $\pm$ 1,5 \\
      & padrão  & 83,3 $\pm$ 0,5 & 88,6 $\pm$ 0,4 & 90,8 $\pm$ 0,1 \\
      & forte   & \textbf{86,8 $\pm$ 0,6} & \textbf{91,5 $\pm$ 0,1} & \textbf{93,3 $\pm$ 0,1} \\
    \multirow{3}{*}{DeiT-Tiny}
      & nenhuma & 52,8 $\pm$ 0,7 & 58,8 $\pm$ 0,5 & 60,9 $\pm$ 2,2 \\
      & padrão  & 59,2 $\pm$ 1,3 & 61,6 $\pm$ 6,6 & 70,2 $\pm$ 3,5 \\
      & forte   & 49,6 $\pm$ 4,0 & 57,8 $\pm$ 8,0 & 68,7 $\pm$ 10,4 \\
    \midrule
    \multicolumn{5}{l}{\textit{D-Fire}} \\
    \multirow{3}{*}{Efficient-CapsNet}
      & nenhuma & 87,3 $\pm$ 0,5 & 89,2 $\pm$ 0,2 & 90,2 $\pm$ 0,5 \\
      & padrão  & \textbf{89,0 $\pm$ 0,5} & 92,8 $\pm$ 0,4 & 93,6 $\pm$ 0,4 \\
      & forte   & 86,4 $\pm$ 0,0 & 90,3 $\pm$ 0,4 & 91,9 $\pm$ 0,3 \\
    \multirow{3}{*}{ResNet-18}
      & nenhuma & 84,4 $\pm$ 1,7 & 88,7 $\pm$ 0,6 & 91,2 $\pm$ 0,4 \\
      & padrão  & 88,8 $\pm$ 0,9 & 92,7 $\pm$ 0,3 & 94,1 $\pm$ 0,5 \\
      & forte   & \textbf{90,0 $\pm$ 0,7} & \textbf{93,8 $\pm$ 0,2} & \textbf{94,7 $\pm$ 0,2} \\
    \multirow{3}{*}{DeiT-Tiny}
      & nenhuma & 82,2 $\pm$ 0,7 & 84,0 $\pm$ 2,5 & 84,6 $\pm$ 2,9 \\
      & padrão  & 76,0 $\pm$ 7,4 & 78,9 $\pm$ 8,8 & 81,0 $\pm$ 5,4 \\
      & forte   & 73,9 $\pm$ 2,0 & 75,1 $\pm$ 5,3 & 77,5 $\pm$ 8,1 \\
    \bottomrule
  \end{tabular}
\end{table}

No CIFAR-10 a ResNet-18 vence em todas as nove combinações de fração e
augmentação, com margem entre $4{,}4$ e $10{,}5$ pontos percentuais (pp) sobre o
Efficient-CapsNet. O melhor resultado capsular ($84{,}8\%$, augmentação padrão,
$f = 1{,}00$) fica $8{,}5$~pp abaixo do melhor resultado convolucional
($93{,}3\%$). O DeiT-Tiny é o pior dos três por larga margem, o que é esperado:
transformadores de visão treinados do zero, sem pré-treino e com 45.000 imagens
de $32\times32$, operam muito abaixo do seu regime.

No D-Fire o quadro é diferente. A tarefa é binária e mais fácil, e o
Efficient-CapsNet fica a $1{,}1$~pp da ResNet-18 na melhor configuração de cada
um ($93{,}6\%$ contra $94{,}7\%$). Em duas células --- augmentação nenhuma com
$f = 0{,}33$ e $f = 0{,}66$, e augmentação padrão com $f = 0{,}33$ --- o modelo
capsular tem a maior média, mas nenhuma dessas vantagens excede o dobro do
ruído entre sementes.

Duas observações estruturais valem o registro. A primeira: a ResNet-18 é a
única arquitetura que melhora ao subir de augmentação padrão para forte.
Em $f = 1{,}00$, a passagem custa $1{,}8$~pp ao Efficient-CapsNet no CIFAR-10 e
$1{,}7$~pp no D-Fire, e custa $1{,}6$ e $3{,}5$~pp ao DeiT-Tiny, enquanto dá
$2{,}5$ e $0{,}6$~pp à ResNet-18. A segunda: o desvio padrão entre sementes do
Efficient-CapsNet é a menor de todas em quase toda a tabela, ponto retomado na
Seção~\ref{sec:res:seeds}.

% ---------------------------------------------------------------------
\section{Robustez sob transformações não vistas}
\label{sec:res:robust}

A Tabela~\ref{tab:res-robust} reporta a condição $f = 1{,}00$ sob as duas
sondas de deslocamento. Além da acurácia absoluta, listamos $R_c$ e a
\emph{habilidade retida}
\begin{equation}
  S_c = \frac{A^{c} - A^{\text{acaso}}}{A^{\text{clean}} - A^{\text{acaso}}},
  \label{eq:skill}
\end{equation}
que desconta o acerto trivial antes de formar a razão. A diferença entre as
duas colunas é instrutiva: $R_c$ dá crédito por acertar ao acaso, $S_c$ não.

\begin{table}[t]
  \centering
  \caption{Desempenho sob transformações não vistas em $f = 1{,}00$, sem
  augmentação em teste. $A^{c}$ é a acurácia absoluta (\%), $R_c$ a razão de
  retenção (\%) e $S_c$ a habilidade retida (\%) da Equação~\eqref{eq:skill}.
  Acaso: $10\%$ no CIFAR-10, $50\%$ no D-Fire.}
  \label{tab:res-robust}
  \small
  \begin{tabular}{llcccccc}
    \toprule
    \multirow{2}{*}{Modelo} & \multirow{2}{*}{Aug.}
      & \multicolumn{3}{c}{Rotação grande} & \multicolumn{3}{c}{Todas as transformações} \\
    \cmidrule(lr){3-5}\cmidrule(lr){6-8}
     & & $A^{c}$ & $R_c$ & $S_c$ & $A^{c}$ & $R_c$ & $S_c$ \\
    \midrule
    \multicolumn{8}{l}{\textit{CIFAR-10}} \\
    \multirow{3}{*}{Efficient-CapsNet}
      & nenhuma & 38,1 & 51,3 & 43,7 & 12,6 & 17,0 & 4,1 \\
      & padrão  & 52,9 & 62,4 & 57,4 & 12,8 & 15,1 & 3,8 \\
      & forte   & 51,4 & 61,9 & 56,6 & 13,8 & 16,7 & 5,2 \\
    \multirow{3}{*}{ResNet-18}
      & nenhuma & 36,4 & 46,2 & 38,4 & 11,6 & 14,7 & 2,3 \\
      & padrão  & 57,8 & 63,6 & 59,1 & 12,3 & 13,6 & 2,9 \\
      & forte   & \textbf{60,0} & 64,4 & 60,1 & 13,4 & 14,4 & 4,1 \\
    \multirow{3}{*}{DeiT-Tiny}
      & nenhuma & 40,0 & 65,8 & 59,0 & \textbf{16,4} & 26,9 & 12,6 \\
      & padrão  & 52,8 & 75,1 & 71,0 & 15,9 & 22,6 & 9,8 \\
      & forte   & 52,0 & \textbf{75,7} & \textbf{71,6} & \textbf{16,8} & \textbf{24,5} & \textbf{11,6} \\
    \midrule
    \multicolumn{8}{l}{\textit{D-Fire}} \\
    \multirow{3}{*}{Efficient-CapsNet}
      & nenhuma & 61,4 & 68,1 & 28,4 & 53,1 & 58,9 & 7,7 \\
      & padrão  & 77,6 & 82,9 & 63,2 & 56,1 & 59,9 & 13,9 \\
      & forte   & 80,0 & 87,0 & 71,5 & 60,5 & 65,8 & 25,0 \\
    \multirow{3}{*}{ResNet-18}
      & nenhuma & 63,4 & 69,5 & 32,4 & 52,5 & 57,5 & 6,0 \\
      & padrão  & 78,7 & 83,7 & 65,2 & 54,9 & 58,4 & 11,2 \\
      & forte   & \textbf{80,9} & 85,4 & 69,1 & \textbf{60,7} & 64,1 & 23,9 \\
    \multirow{3}{*}{DeiT-Tiny}
      & nenhuma & 63,5 & 75,0 & 38,9 & 52,9 & 62,5 & 8,4 \\
      & padrão  & 78,4 & 96,9 & 91,8 & 53,7 & 66,4 & 12,1 \\
      & forte   & 76,0 & \textbf{98,0} & \textbf{94,4} & 59,9 & \textbf{77,2} & \textbf{35,9} \\
    \bottomrule
  \end{tabular}
\end{table}

\subsection{A sonda de rotação}

A rotação grande é o teste direto de H1, e o resultado é desfavorável à
hipótese. No CIFAR-10 com augmentação padrão ou forte, a ResNet-18 tem a maior
acurácia absoluta sob rotação --- $57{,}8\%$ e $60{,}0\%$ contra $52{,}9\%$ e
$51{,}4\%$ do Efficient-CapsNet. A vantagem convolucional sob rotação
($+8{,}6$~pp em $f = 1{,}00$, augmentação forte) é maior, em valor absoluto, do
que qualquer vantagem capsular observada em toda a grade.

Apenas no regime sem augmentação nenhuma o Efficient-CapsNet fica à frente sob
rotação, e por pouco: $38{,}1\%$ contra $36{,}4\%$ da ResNet-18. Em teste
pareado por semente sobre todas as 27 células do CIFAR-10, a diferença de
retenção entre Efficient-CapsNet e ResNet-18 sob rotação é de $+1{,}1$~pp com
$p = 0{,}13$ --- indistinguível de zero. No D-Fire a mesma comparação dá
$-1{,}5$~pp com $p = 0{,}013$, ou seja, uma desvantagem capsular
estatisticamente detectável. Contra o DeiT-Tiny a diferença de retenção é de
$-15{,}1$~pp no CIFAR-10 e $-11{,}7$~pp no D-Fire, ambas com $p < 10^{-12}$.

A Figura~\ref{fig:res-dumbbell} torna o padrão visível: as barras do
Efficient-CapsNet e da ResNet-18 partem de pontos diferentes e caem
aproximadamente na mesma proporção, sem que a distância entre elas se feche.

\begin{figure}[t]
  \centering
  \includegraphics[width=\textwidth]{fig01_shift_dumbbell_cifar_10.png}
  \caption{CIFAR-10: acurácia limpa (círculo aberto) e acurácia sob
  deslocamento (marcador preenchido) para cada combinação de arquitetura,
  fração de dados e regime de augmentação. O comprimento da barra é a
  degradação absoluta.}
  \label{fig:res-dumbbell}
\end{figure}

\subsection{A composição de todas as transformações é um piso, não uma medida}
\label{sec:res:floor}

Sob a composição das seis corrupções, os três modelos deixam de funcionar. No
CIFAR-10 a acurácia cai para a faixa de $11{,}1\%$ a $19{,}1\%$, contra
$10{,}0\%$ de acaso, em todas as 54 células. No D-Fire os valores ficam entre
$-2{,}2$ e $+7{,}8$~pp da linha de base da classe majoritária ($53{,}4\%$); sem
augmentação de treino, 14 das 18 células ficam \emph{abaixo} dela, isto é,
perde para um classificador que ignora a imagem
(Figura~\ref{fig:res-baselines}).

Isso muda o que se pode concluir a partir dessa condição. Comparar retenção
entre modelos que estão todos no nível do acaso compara padrões de colapso, não
robustez. As diferenças de $R_c$ que aparecem ali --- $26{,}9\%$ para o
DeiT-Tiny contra $14{,}7\%$ para a ResNet-18, por exemplo --- correspondem, em
acurácia absoluta, a $16{,}4\%$ contra $11{,}6\%$: quatro pontos e meio acima e
um ponto e meio acima do acaso, respectivamente. Nenhum dos dois é utilizável.

\begin{figure}[t]
  \centering
  \includegraphics[width=\textwidth]{fig11_cross_dataset_vs_baseline.png}
  \caption{Os dois conjuntos contra suas próprias linhas de base triviais.
  Barras abaixo da linha pontilhada perdem para um classificador que ignora a
  imagem. Sob todas as transformações, nenhuma arquitetura se sustenta em
  nenhum dos conjuntos.}
  \label{fig:res-baselines}
\end{figure}

A recomendação que decorre daí é metodológica e vale além deste trabalho:
composições de corrupções só informam enquanto os modelos comparados
permanecerem acima do acaso por margem apreciável. Abaixo disso, a métrica
mede outra coisa. A Seção~\ref{sec:res:collapse} identifica exatamente o que.

\subsection{Por que a retenção não pode ser a métrica principal}
\label{sec:res:confound}

O DeiT-Tiny tem a maior razão de retenção em praticamente toda a
Tabela~\ref{tab:res-robust} e a menor acurácia limpa em toda a
Tabela~\ref{tab:res-clean}. Não é coincidência. A Figura~\ref{fig:res-confound}
mostra a correlação entre acurácia limpa e retenção no CIFAR-10:
$r = -0{,}44$ sob rotação e $r = -0{,}91$ sob todas as transformações. As
inclinações ajustadas são de $-34$ e $-48$~pp de retenção por unidade de
acurácia limpa.

\begin{figure}[t]
  \centering
  \includegraphics[width=\textwidth]{fig03_retention_confound_cifar_10.png}
  \caption{CIFAR-10: razão de retenção contra acurácia limpa. A relação
  negativa significa que a retenção, isolada, recompensa modelos que já eram
  ruins em distribuição.}
  \label{fig:res-confound}
\end{figure}

O mecanismo é aritmético. A acurácia sob deslocamento tem piso no acaso,
enquanto a acurácia limpa não tem teto próximo; um modelo que sai de $60\%$ e
cai para $40\%$ retém $67\%$, e um que sai de $93\%$ e cai para $60\%$ retém
$64\%$ --- embora o segundo esteja $20$~pp à frente onde importa. A
Figura~\ref{fig:res-frontier} mostra a mesma informação no plano
(acurácia limpa, acurácia deslocada): sob rotação a correlação entre os dois
eixos é $+0{,}68$, e sob todas as transformações ela \emph{inverte} de sinal
($-0{,}74$), que é a assinatura do regime de colapso descrito acima.

\begin{figure}[t]
  \centering
  \includegraphics[width=\textwidth]{fig02_robustness_frontier_cifar_10.png}
  \caption{CIFAR-10: acurácia sob deslocamento contra acurácia limpa. A
  diagonal tracejada marca invariância perfeita. Separação vertical entre cores
  na mesma abscissa seria o efeito arquitetural puro; ela é pequena sob rotação
  e some sob todas as transformações.}
  \label{fig:res-frontier}
\end{figure}

Por isso a Tabela~\ref{tab:res-robust} reporta $A^c$ em primeiro lugar, e por
isso H1 é avaliada sobre acurácia absoluta e sobre retenção pareada por
semente, não sobre médias de $R_c$.

% ---------------------------------------------------------------------
\section{Colapso de predição e detecção de erro}
\label{sec:res:collapse}

Dois diagnósticos separam ``degradação graciosa'' de ``falha silenciosa'': a
fração de predições que caem na classe mais predita (\emph{top-pred share}) e a
AUROC da confiança do modelo como detector dos próprios erros. A
Tabela~\ref{tab:res-auroc} traz os dois.

\begin{table}[t]
  \centering
  \caption{AUROC de detecção de erro (a confiança separa acertos de erros?) e
  fração de predições na classe mais predita, médias sobre as três frações de
  dados, sem augmentação em teste. AUROC de $0{,}50$ é o acaso; \emph{share} de
  $0{,}10$ (CIFAR-10) e $0{,}53$ (D-Fire) corresponde a predições
  balanceadas.}
  \label{tab:res-auroc}
  \footnotesize
  \setlength{\tabcolsep}{4pt}
  \begin{tabular}{lccccccccc}
    \toprule
    \multirow{2}{*}{Modelo}
      & \multicolumn{3}{c}{Limpo} & \multicolumn{3}{c}{Rotação grande} & \multicolumn{3}{c}{Todas} \\
    \cmidrule(lr){2-4}\cmidrule(lr){5-7}\cmidrule(lr){8-10}
     & nenh. & padr. & forte & nenh. & padr. & forte & nenh. & padr. & forte \\
    \midrule
    \multicolumn{10}{l}{\textit{CIFAR-10 --- AUROC}} \\
    Efficient-CapsNet & 0,799 & 0,856 & 0,848 & 0,707 & 0,762 & 0,766 & 0,533 & 0,547 & 0,548 \\
    ResNet-18         & 0,788 & 0,851 & 0,867 & 0,697 & 0,790 & 0,804 & 0,540 & 0,527 & 0,540 \\
    DeiT-Tiny         & 0,716 & 0,736 & 0,720 & 0,677 & 0,722 & 0,713 & 0,560 & 0,587 & 0,618 \\
    \addlinespace
    \multicolumn{10}{l}{\textit{CIFAR-10 --- fração na classe mais predita}} \\
    Efficient-CapsNet & 0,12 & 0,11 & 0,11 & 0,24 & 0,18 & 0,20 & 0,39 & 0,48 & 0,43 \\
    ResNet-18         & 0,13 & 0,11 & 0,11 & 0,28 & 0,18 & 0,18 & 0,55 & 0,70 & 0,76 \\
    DeiT-Tiny         & 0,13 & 0,12 & 0,13 & 0,21 & 0,15 & 0,16 & 0,38 & 0,48 & 0,38 \\
    \midrule
    \multicolumn{10}{l}{\textit{D-Fire --- AUROC}} \\
    Efficient-CapsNet & 0,837 & 0,838 & 0,836 & 0,581 & 0,696 & 0,728 & 0,515 & 0,558 & 0,595 \\
    ResNet-18         & 0,826 & 0,845 & 0,881 & 0,610 & 0,736 & 0,745 & 0,532 & 0,562 & 0,601 \\
    DeiT-Tiny         & 0,769 & 0,755 & 0,731 & 0,627 & 0,727 & 0,726 & 0,525 & 0,560 & 0,566 \\
    \addlinespace
    \multicolumn{10}{l}{\textit{D-Fire --- fração na classe mais predita}} \\
    Efficient-CapsNet & 0,56 & 0,54 & 0,55 & 0,78 & 0,64 & 0,62 & 0,75 & 0,60 & 0,58 \\
    ResNet-18         & 0,53 & 0,53 & 0,53 & 0,70 & 0,61 & 0,62 & 0,65 & 0,58 & 0,62 \\
    DeiT-Tiny         & 0,55 & 0,54 & 0,54 & 0,55 & 0,54 & 0,54 & 0,70 & 0,63 & 0,63 \\
    \bottomrule
  \end{tabular}
\end{table}

Em dados limpos, o Efficient-CapsNet tem a melhor detecção de erro do CIFAR-10
sob augmentação padrão ($0{,}856$ contra $0{,}851$ da ResNet-18) e a melhor do
D-Fire no regime sem augmentação ($0{,}837$ contra $0{,}826$), ficando a menos
de um centésimo da ResNet-18 sob augmentação padrão. A norma do vetor de cápsula, portanto,
funciona como escore de confiança pelo menos tão bem quanto o \emph{softmax},
o que é um resultado positivo para o modelo capsular e independente da
acurácia.

Sob deslocamento essa propriedade não sobrevive. Sob rotação a AUROC capsular
cai para $0{,}707$--$0{,}766$ no CIFAR-10 e para $0{,}581$ no D-Fire sem
augmentação --- neste último caso, com $77{,}6\%$ das predições em uma única
classe. Ou seja: o modelo rotacionado no D-Fire responde quase sempre a mesma
coisa e não sabe que está errando. Sob todas as transformações, a AUROC dos
três modelos fica entre $0{,}515$ e $0{,}618$, praticamente o acaso
(Figura~\ref{fig:res-auroc}). A falha, nessa condição, é silenciosa em todas as
arquiteturas.

\begin{figure}[t]
  \centering
  \includegraphics[width=\textwidth]{fig07_error_detection_auroc_cifar_10.png}
  \caption{CIFAR-10: AUROC da confiança como detector dos próprios erros. A
  linha vermelha marca o acaso. Sob a composição de transformações, nenhum
  modelo distingue acerto de erro.}
  \label{fig:res-auroc}
\end{figure}

A fração na classe mais predita revela a diferença de estilo de falha. Sob
todas as transformações no CIFAR-10, a ResNet-18 colapsa: $76\%$ das predições
em uma única classe no regime de augmentação forte, contra $43\%$ do
Efficient-CapsNet e $38\%$ do DeiT-Tiny (Figura~\ref{fig:res-collapse}). A
Figura~\ref{fig:res-perclass} mostra a forma desse colapso por classe: a
ResNet-18 corrompida acerta $42\%$ de \emph{airplane} e $55\%$ de \emph{bird} e
fica em $0$--$3\%$ em sete das dez classes. O Efficient-CapsNet degrada de
modo mais distribuído, sem por isso acertar mais.

\begin{figure}[t]
  \centering
  \includegraphics[width=\textwidth]{fig08_prediction_collapse_cifar_10.png}
  \caption{CIFAR-10: fração das predições concentradas na classe mais predita.
  Valores próximos de $1/K$ indicam predições ainda distribuídas; valores altos
  indicam colapso.}
  \label{fig:res-collapse}
\end{figure}

\begin{figure}[t]
  \centering
  \includegraphics[width=\textwidth]{fig09_per_class_accuracy_cifar_10.png}
  \caption{CIFAR-10: acurácia por classe e condição de teste, sem augmentação
  de treino, $f = 1{,}00$, agregada sobre sementes.}
  \label{fig:res-perclass}
\end{figure}

Um detalhe da Figura~\ref{fig:res-perclass} merece nota. A classe \emph{frog} é
a única que o Efficient-CapsNet mantém sob rotação grande: $85\%$ limpo e
$85\%$ rotacionado, retenção integral. A ResNet-18 mostra o mesmo padrão nessa
classe ($86\% \to 80\%$). A explicação mais econômica não é arquitetural, e sim
do dado: sapos aparecem em poses e orientações variadas e sem eixo vertical
dominante, ao contrário de cavalos, navios e caminhões. Onde o prior de
orientação é fraco, a rotação não penaliza --- em qualquer arquitetura.

% ---------------------------------------------------------------------
\section{Eficiência de dados}
\label{sec:res:data}

A Tabela~\ref{tab:res-dataeff} traz a estatística $\Delta_f$ da
Equação~\eqref{eq:data-eff}, a queda de acurácia limpa entre $f = 1{,}00$ e
$f = 0{,}33$.

\begin{table}[t]
  \centering
  \caption{Eficiência de dados: $\Delta_f = A_{f=1{,}00} - A_{f=0{,}33}$ em pp
  de acurácia limpa. Valores menores indicam menor dependência do volume de
  dados.}
  \label{tab:res-dataeff}
  \small
  \begin{tabular}{llccc c}
    \toprule
    \multirow{2}{*}{Conjunto} & \multirow{2}{*}{Modelo} & \multicolumn{3}{c}{Augmentação} & \multirow{2}{*}{Média} \\
    \cmidrule(lr){3-5}
     & & nenhuma & padrão & forte & \\
    \midrule
    \multirow{3}{*}{CIFAR-10}
      & Efficient-CapsNet & 10,4 & \phantom{0}6,7 & \phantom{0}6,5 & \textbf{7,9} \\
      & ResNet-18         & 12,6 & \phantom{0}7,5 & \phantom{0}6,5 & 8,9 \\
      & DeiT-Tiny         & \phantom{0}8,1 & 11,0 & 19,0 & 12,7 \\
    \midrule
    \multirow{3}{*}{D-Fire}
      & Efficient-CapsNet & \phantom{0}2,9 & \phantom{0}4,6 & \phantom{0}5,4 & 4,3 \\
      & ResNet-18         & \phantom{0}6,8 & \phantom{0}5,3 & \phantom{0}4,7 & 5,6 \\
      & DeiT-Tiny         & \phantom{0}2,4 & \phantom{0}5,0 & \phantom{0}3,6 & \textbf{3,7} \\
    \bottomrule
  \end{tabular}
\end{table}

O Efficient-CapsNet tem o menor $\Delta_f$ médio no CIFAR-10 ($7{,}9$ contra
$8{,}9$~pp da ResNet-18) e fica $1{,}3$~pp à frente dela no D-Fire. A direção é
a prevista por H2, mas a magnitude não sustenta a hipótese: no teste pareado
por semente e regime de augmentação, a diferença Efficient-CapsNet menos
ResNet-18 é de $-1{,}0$~pp com $p = 0{,}080$ no CIFAR-10 e $-1{,}3$~pp com
$p = 0{,}139$ no D-Fire, ambas com $n = 9$. Nenhuma alcança significância
convencional, e o poder com três sementes é baixo demais para que a ausência de
significância seja informativa em qualquer direção.

Além disso, o ordenamento absoluto não se altera em nenhuma fração: a
ResNet-18 em $f = 0{,}33$ já supera o Efficient-CapsNet em $f = 1{,}00$ no
CIFAR-10 sob augmentação forte ($86{,}8\%$ contra $83{,}0\%$). Uma eventual
inclinação mais suave da curva capsular não se converte em vantagem prática
dentro da faixa investigada.

A Figura~\ref{fig:res-scaling} mostra também um efeito que $\Delta_f$ não
captura: sob rotação grande, mais dados continuam ajudando a ResNet-18
($+4{,}9$~pp de $f = 0{,}33$ a $f = 1{,}00$ com augmentação forte), enquanto
sob todas as transformações as curvas são planas ou levemente decrescentes para
os três modelos. Dados adicionais compram acurácia em distribuição e sob
deslocamento moderado; não compram nada no regime de colapso.

\begin{figure}[t]
  \centering
  \includegraphics[width=\textwidth]{fig05_data_scaling_cifar_10.png}
  \caption{CIFAR-10: acurácia contra fração de dados de treino, por regime de
  augmentação (linhas) e condição de teste (colunas). Barras de erro são o
  desvio padrão entre três sementes.}
  \label{fig:res-scaling}
\end{figure}

% ---------------------------------------------------------------------
\section{Interação com a augmentação}
\label{sec:res:aug}

H3 previa que uma eventual vantagem capsular encolheria ao nível do ruído
quando as linhas de base recebessem augmentação forte. O resultado observado é
mais forte que a hipótese e aponta na direção contrária: não há vantagem
capsular a encolher, e a \emph{desvantagem} capsular \emph{cresce} com a
augmentação.

A Figura~\ref{fig:res-ladder} quantifica o quanto cada arquitetura converte
augmentação em desempenho. Subindo de ``nenhuma'' a ``forte'' no CIFAR-10, a
ResNet-18 ganha $17$~pp em dados limpos e $22$~pp sob rotação; o
Efficient-CapsNet ganha $11$ e $12$~pp; o DeiT-Tiny ganha $1$ e $9$~pp. A
ResNet-18 é, das três, a que melhor aproveita o regime.

\begin{figure}[t]
  \centering
  \includegraphics[width=\textwidth]{fig04_augmentation_ladder_cifar_10.png}
  \caption{CIFAR-10: efeito da intensidade da augmentação de treino sobre cada
  condição de teste, com média sobre as frações de dados. As anotações dão o
  ganho total de ``nenhuma'' a ``forte''.}
  \label{fig:res-ladder}
\end{figure}

A consequência aparece na Figura~\ref{fig:res-heatmap}, que mostra a diferença
entre o Efficient-CapsNet e a melhor linha de base em cada uma das 54
condições do CIFAR-10. Sem augmentação e com $f = 0{,}33$, a diferença é de
$-2{,}3$~pp em dados limpos e chega a $+2{,}3$~pp sob rotação, ambas dentro do
dobro do ruído entre sementes. Com augmentação forte e $f = 1{,}00$, ela é de
$-10{,}3$~pp em dados limpos e $-8{,}7$~pp sob rotação. Em média sobre as
frações, a desvantagem capsular em dados limpos passa de $-2{,}5$~pp (sem
augmentação) para $-8{,}8$ (padrão) e $-13{,}2$~pp (forte).

\begin{figure}[t]
  \centering
  \includegraphics[width=\textwidth]{fig06_advantage_heatmap_cifar_10.png}
  \caption{CIFAR-10: diferença de acurácia entre o Efficient-CapsNet e a melhor
  linha de base em cada condição. Azul indica vantagem capsular, vermelho
  desvantagem; o símbolo $^\circ$ marca células dentro do dobro do ruído entre
  sementes.}
  \label{fig:res-heatmap}
\end{figure}

A interpretação natural é que a invariância que as cápsulas ofereceriam por
construção pode ser adquirida pelas linhas de base por meio de dados, e que
adquiri-la por dados é mais barato e mais eficaz. O resultado é consistente com
a leitura crítica da literatura: o que se credita ao roteamento talvez seja, em
boa parte, atribuível ao objetivo de treino e ao regime de augmentação
(Seção~\ref{sec:protocol:threats}).

Vale registrar o contraponto: nem o Efficient-CapsNet nem o DeiT-Tiny se
beneficiam de augmentação forte em dados limpos; ambos pioram em relação à
augmentação padrão, nos dois conjuntos. Isso
sugere que a perda de margem com reconstrução, ou a geometria das cápsulas
primárias, tolera mal o RandAugment e o Random Erasing --- e que a comparação
sob ``augmentação forte'' talvez não esteja comparando cada arquitetura em seu
melhor. É uma limitação do delineamento, não um resultado sobre cápsulas.

% ---------------------------------------------------------------------
\section{Variabilidade entre sementes}
\label{sec:res:seeds}

A Tabela~\ref{tab:res-seeds} reporta a amplitude entre as três sementes, média
sobre todas as condições de avaliação.

\begin{table}[t]
  \centering
  \caption{Amplitude (máximo menos mínimo) da acurácia entre três sementes, em
  pp, média sobre todas as condições de avaliação.}
  \label{tab:res-seeds}
  \small
  \begin{tabular}{llcccc}
    \toprule
    \multirow{2}{*}{Conjunto} & \multirow{2}{*}{Modelo} & \multicolumn{3}{c}{Augmentação} & \multirow{2}{*}{Média} \\
    \cmidrule(lr){3-5}
     & & nenhuma & padrão & forte & \\
    \midrule
    \multirow{3}{*}{CIFAR-10}
      & Efficient-CapsNet & 0,93 & 0,93 & 0,86 & \textbf{0,91} \\
      & ResNet-18         & 2,01 & 1,17 & 1,06 & 1,42 \\
      & DeiT-Tiny         & 1,72 & 3,39 & 6,75 & 3,95 \\
    \midrule
    \multirow{3}{*}{D-Fire}
      & Efficient-CapsNet & 1,44 & 1,62 & 1,38 & \textbf{1,48} \\
      & ResNet-18         & 1,84 & 1,26 & 1,79 & 1,63 \\
      & DeiT-Tiny         & 2,14 & 7,57 & 7,70 & 5,80 \\
    \bottomrule
  \end{tabular}
\end{table}

Este é o resultado mais nítido a favor do modelo capsular em toda a grade, e
não é sobre robustez nem sobre acurácia. O Efficient-CapsNet é a arquitetura
mais estável entre sementes nos dois conjuntos, com amplitude média de
$0{,}91$~pp no CIFAR-10 --- menos de dois terços da ResNet-18 e menos de um
quarto do DeiT-Tiny. A estabilidade é quase insensível ao regime de
augmentação, enquanto a do DeiT-Tiny se deteriora fortemente com ele (de
$1{,}7$ para $6{,}8$~pp no CIFAR-10).

A amplitude individual mais extrema da grade é de $19{,}4$~pp, no DeiT-Tiny com
augmentação forte e $f = 1{,}00$ no CIFAR-10 --- uma única semente separa
$57\%$ de $76\%$ de acurácia. Os doze piores casos de dispersão da
grade pertencem todos ao DeiT-Tiny, oito deles no D-Fire. Qualquer número reportado para essa arquitetura neste
trabalho deve ser lido com essa dispersão em mente, e a
Figura~\ref{fig:res-seeds} torna explícito que várias comparações envolvendo
o DeiT-Tiny não sobrevivem ao ruído.

\begin{figure}[t]
  \centering
  \includegraphics[width=\textwidth]{fig10_seed_consistency_cifar_10.png}
  \caption{CIFAR-10, $f = 1{,}00$: cada semente desenhada individualmente,
  ligando a acurácia limpa à acurácia sob todas as transformações.}
  \label{fig:res-seeds}
\end{figure}

Uma leitura possível é que a perda de margem com reconstrução, por regularizar
mais fortemente, produz um objetivo com paisagem menos sensível à
inicialização. A hipótese não é testável dentro deste delineamento, que
confunde arquitetura e função objetivo, mas é a explicação candidata a ser
investigada.

% ---------------------------------------------------------------------
\section{Contagem de vitórias}
\label{sec:res:wins}

Somando as 108 condições de avaliação (2 conjuntos $\times$ 3 augmentações
$\times$ 3 frações $\times$ 2 augmentações em teste $\times$ 3 transformações),
a ResNet-18 vence em 57, o DeiT-Tiny em 33 e o Efficient-CapsNet em 18. Das 18
vitórias capsulares, \emph{uma única} excede o dobro do ruído entre sementes:
D-Fire, augmentação forte, $f = 1{,}00$, teste com augmentação forte e todas as
transformações, $59{,}2\%$ contra $57{,}7\%$ da ResNet-18 --- e essa é
justamente uma condição em que ambos os modelos estão a poucos pontos da linha
de base da classe majoritária.

A distribuição das vitórias também localiza os regimes. As 33 vitórias do
DeiT-Tiny concentram-se sob deslocamento (18 delas na composição de todas as
transformações no CIFAR-10, onde todos estão perto do acaso). As 57 vitórias da
ResNet-18 concentram-se em dados limpos e sob rotação, isto é, nos regimes em
que os números significam alguma coisa.

% ---------------------------------------------------------------------
\section{Veredito sobre as hipóteses}
\label{sec:res:verdict}

\begin{table}[t]
  \centering
  \caption{Resumo da avaliação das hipóteses da
  Seção~\ref{sec:protocol:hypotheses}.}
  \label{tab:res-verdict}
  \small
  \begin{tabular}{l p{0.68\textwidth}}
    \toprule
    Hipótese & Veredito \\
    \midrule
    H1 & \textbf{Rejeitada.} Sob rotação grande, a retenção capsular não supera
    a da ResNet-18 no CIFAR-10 ($+1{,}1$~pp, $p = 0{,}13$) e é inferior no
    D-Fire ($-1{,}5$~pp, $p = 0{,}013$). Em acurácia absoluta, a ResNet-18 está
    à frente em seis das nove condições do CIFAR-10, incluindo as seis com
    augmentação. \\
    \addlinespace
    H2 & \textbf{Não sustentada.} A direção prevista aparece ($\Delta_f$ menor
    para as cápsulas em ambos os conjuntos), mas a diferença é de cerca de
    $1$~pp e não atinge significância com três sementes
    ($p = 0{,}080$ e $p = 0{,}139$). O ordenamento absoluto não se inverte em
    nenhuma fração. \\
    \addlinespace
    H3 & \textbf{Rejeitada, com inversão de sinal.} Não há vantagem capsular a
    encolher. A desvantagem em dados limpos no CIFAR-10 \emph{cresce} de
    $-2{,}5$~pp sem augmentação para $-13{,}2$~pp sob augmentação forte. \\
    \bottomrule
  \end{tabular}
\end{table}

O achado central é negativo e deve ser enunciado sem atenuação: nesta grade,
substituir camadas convolucionais por cápsulas com roteamento por
autoatenção não compra robustez a transformações não vistas, nem eficiência de
dados suficiente para compensar a perda de acurácia em distribuição. Onde o
modelo capsular aparece à frente, ou a margem está dentro do ruído entre
sementes, ou os modelos comparados estão no nível do acaso.

Três achados positivos sobrevivem à mesma análise, e nenhum deles é sobre
equivariância:

\begin{description}
\item[Estabilidade.] O Efficient-CapsNet é a arquitetura menos sensível à
semente nos dois conjuntos, com amplitude média de $0{,}91$~pp no CIFAR-10
contra $1{,}42$ da ResNet-18 e $3{,}95$ do DeiT-Tiny
(Seção~\ref{sec:res:seeds}).

\item[Confiança calibrada em distribuição.] A norma do vetor de cápsula é um
escore de detecção de erro competitivo ou superior ao \emph{softmax} em dados
limpos, com AUROC de até $0{,}856$ no CIFAR-10 e $0{,}838$ no D-Fire
(Seção~\ref{sec:res:collapse}).

\item[Competitividade na tarefa binária.] No D-Fire, a diferença para a melhor
linha de base é de $0{,}8$--$3{,}4$~pp conforme o regime, contra
$2{,}5$--$13{,}2$~pp no CIFAR-10. A penalidade capsular é fortemente dependente
da tarefa, o que sugere que conclusões obtidas apenas no CIFAR-10 não
transferem.
\end{description}

% ---------------------------------------------------------------------
\section{Ameaças à validade destes resultados}
\label{sec:res:threats}

Além das ameaças já declaradas na Seção~\ref{sec:protocol:threats}, a análise
deste capítulo acrescenta quatro.

\begin{description}

\item[A composição de corrupções satura.] Sob todas as transformações, os três
modelos ficam no nível do acaso (CIFAR-10) ou da classe majoritária (D-Fire).
Comparações nessa condição não separam robustez; separam estilo de colapso.
Trinta e seis das 108 condições da Seção~\ref{sec:res:wins} pertencem a esse
regime e foram desconsideradas nas conclusões substantivas.

\item[Uma única magnitude de rotação.] A sonda amostra
$\theta \sim \mathcal{U}(-90^{\circ}, 90^{\circ})$ uma vez por imagem. Uma
varredura determinística em $\theta$ produziria uma curva, e é possível que as
cápsulas se saiam melhor em ângulos pequenos e intermediários, onde a
degradação convolucional ainda é suave. O resultado aqui reportado vale para o
regime agregado, não para cada ângulo.

\item[Poder estatístico.] Com três sementes, diferenças abaixo de
aproximadamente $2$~pp são indistinguíveis do ruído na maior parte da grade.
Isso protege contra falsos positivos e impede, simetricamente, afirmar que H2
é falsa: só é possível afirmar que ela não foi demonstrada.

\item[A augmentação forte não é neutra.] Das três arquiteturas, só a ResNet-18
tem a augmentação forte como melhor regime; o Efficient-CapsNet tem a padrão,
nos dois conjuntos. Comparar as três sob augmentação forte é comparar a
ResNet-18 em seu melhor com o Efficient-CapsNet fora dele. A comparação mais favorável ao modelo capsular --- melhor regime
contra melhor regime --- ainda o deixa $8{,}5$~pp atrás no CIFAR-10 e
$1{,}1$~pp atrás no D-Fire, de modo que a conclusão não se inverte; mas as
magnitudes da Seção~\ref{sec:res:aug} são um teto, não uma estimativa pontual.

\end{description}
 a partir do seu .tex principal.
%
%  Este arquivo NÃO carrega pacote nenhum e NÃO declara \documentclass.
%  Ele usa apenas o que o protocolo-experimental.tex já exige:
%    booktabs, multirow, graphicx  (e, se você usar as figuras, um
%    \graphicspath{{figuras/}} no preâmbulo ou o caminho corrigido nos
%    \includegraphics abaixo).
%
%  Se a sua classe for `article`, troque \chapter -> \section e rebaixe
%  os demais níveis.
% =====================================================================

\chapter{Resultados e Discussão}
\label{ch:results}

Este capítulo reporta a grade completa descrita no
Capítulo~\ref{ch:protocol}: 162 treinos (3 arquiteturas $\times$ 2 conjuntos
$\times$ 3 frações de dados $\times$ 3 regimes de augmentação $\times$ 3
sementes), cada um avaliado sob três condições de teste --- limpo, rotação
grande e composição de todas as transformações não vistas --- com e sem
augmentação forte aplicada em teste. São 324 células de resultado e 972
avaliações individuais por semente.

A leitura que segue está organizada em torno de uma decisão metodológica que
precisa ser explicitada antes dos números: \emph{a razão de retenção
$R_c$ da Equação~\eqref{eq:retention} não é usada como métrica principal}.
A Seção~\ref{sec:res:confound} mostra que, nesta grade, $R_c$ é
sistematicamente anticorrelacionada com a acurácia limpa e portanto premia
modelos fracos. As conclusões substantivas repousam sobre a acurácia absoluta
sob deslocamento, lida contra o acaso e contra a linha de base da classe
majoritária.

% ---------------------------------------------------------------------
\section{Desempenho em distribuição}
\label{sec:res:clean}

A Tabela~\ref{tab:res-clean} traz a acurácia de teste limpa. O ordenamento é
estável em toda a grade e não depende da fração de dados.

\begin{table}[t]
  \centering
  \caption{Acurácia de teste limpa (\%) por arquitetura, fração de dados e
  regime de augmentação. Média $\pm$ desvio padrão sobre três sementes.}
  \label{tab:res-clean}
  \small
  \begin{tabular}{llccc}
    \toprule
    \multirow{2}{*}{Modelo} & \multirow{2}{*}{Aug.} & \multicolumn{3}{c}{Fração de dados $f$} \\
    \cmidrule(lr){3-5}
     & & $0{,}33$ & $0{,}66$ & $1{,}00$ \\
    \midrule
    \multicolumn{5}{l}{\textit{CIFAR-10}} \\
    \multirow{3}{*}{Efficient-CapsNet}
      & nenhuma & 63,8 $\pm$ 0,1 & 70,3 $\pm$ 0,4 & 74,2 $\pm$ 0,1 \\
      & padrão  & 78,0 $\pm$ 0,5 & 82,2 $\pm$ 0,1 & 84,8 $\pm$ 0,4 \\
      & forte   & 76,6 $\pm$ 0,3 & 80,7 $\pm$ 0,1 & 83,0 $\pm$ 0,3 \\
    \multirow{3}{*}{ResNet-18}
      & nenhuma & 66,2 $\pm$ 1,9 & 75,3 $\pm$ 0,9 & 78,8 $\pm$ 1,5 \\
      & padrão  & 83,3 $\pm$ 0,5 & 88,6 $\pm$ 0,4 & 90,8 $\pm$ 0,1 \\
      & forte   & \textbf{86,8 $\pm$ 0,6} & \textbf{91,5 $\pm$ 0,1} & \textbf{93,3 $\pm$ 0,1} \\
    \multirow{3}{*}{DeiT-Tiny}
      & nenhuma & 52,8 $\pm$ 0,7 & 58,8 $\pm$ 0,5 & 60,9 $\pm$ 2,2 \\
      & padrão  & 59,2 $\pm$ 1,3 & 61,6 $\pm$ 6,6 & 70,2 $\pm$ 3,5 \\
      & forte   & 49,6 $\pm$ 4,0 & 57,8 $\pm$ 8,0 & 68,7 $\pm$ 10,4 \\
    \midrule
    \multicolumn{5}{l}{\textit{D-Fire}} \\
    \multirow{3}{*}{Efficient-CapsNet}
      & nenhuma & 87,3 $\pm$ 0,5 & 89,2 $\pm$ 0,2 & 90,2 $\pm$ 0,5 \\
      & padrão  & \textbf{89,0 $\pm$ 0,5} & 92,8 $\pm$ 0,4 & 93,6 $\pm$ 0,4 \\
      & forte   & 86,4 $\pm$ 0,0 & 90,3 $\pm$ 0,4 & 91,9 $\pm$ 0,3 \\
    \multirow{3}{*}{ResNet-18}
      & nenhuma & 84,4 $\pm$ 1,7 & 88,7 $\pm$ 0,6 & 91,2 $\pm$ 0,4 \\
      & padrão  & 88,8 $\pm$ 0,9 & 92,7 $\pm$ 0,3 & 94,1 $\pm$ 0,5 \\
      & forte   & \textbf{90,0 $\pm$ 0,7} & \textbf{93,8 $\pm$ 0,2} & \textbf{94,7 $\pm$ 0,2} \\
    \multirow{3}{*}{DeiT-Tiny}
      & nenhuma & 82,2 $\pm$ 0,7 & 84,0 $\pm$ 2,5 & 84,6 $\pm$ 2,9 \\
      & padrão  & 76,0 $\pm$ 7,4 & 78,9 $\pm$ 8,8 & 81,0 $\pm$ 5,4 \\
      & forte   & 73,9 $\pm$ 2,0 & 75,1 $\pm$ 5,3 & 77,5 $\pm$ 8,1 \\
    \bottomrule
  \end{tabular}
\end{table}

No CIFAR-10 a ResNet-18 vence em todas as nove combinações de fração e
augmentação, com margem entre $4{,}4$ e $10{,}5$ pontos percentuais (pp) sobre o
Efficient-CapsNet. O melhor resultado capsular ($84{,}8\%$, augmentação padrão,
$f = 1{,}00$) fica $8{,}5$~pp abaixo do melhor resultado convolucional
($93{,}3\%$). O DeiT-Tiny é o pior dos três por larga margem, o que é esperado:
transformadores de visão treinados do zero, sem pré-treino e com 45.000 imagens
de $32\times32$, operam muito abaixo do seu regime.

No D-Fire o quadro é diferente. A tarefa é binária e mais fácil, e o
Efficient-CapsNet fica a $1{,}1$~pp da ResNet-18 na melhor configuração de cada
um ($93{,}6\%$ contra $94{,}7\%$). Em duas células --- augmentação nenhuma com
$f = 0{,}33$ e $f = 0{,}66$, e augmentação padrão com $f = 0{,}33$ --- o modelo
capsular tem a maior média, mas nenhuma dessas vantagens excede o dobro do
ruído entre sementes.

Duas observações estruturais valem o registro. A primeira: a ResNet-18 é a
única arquitetura que melhora ao subir de augmentação padrão para forte.
Em $f = 1{,}00$, a passagem custa $1{,}8$~pp ao Efficient-CapsNet no CIFAR-10 e
$1{,}7$~pp no D-Fire, e custa $1{,}6$ e $3{,}5$~pp ao DeiT-Tiny, enquanto dá
$2{,}5$ e $0{,}6$~pp à ResNet-18. A segunda: o desvio padrão entre sementes do
Efficient-CapsNet é a menor de todas em quase toda a tabela, ponto retomado na
Seção~\ref{sec:res:seeds}.

% ---------------------------------------------------------------------
\section{Robustez sob transformações não vistas}
\label{sec:res:robust}

A Tabela~\ref{tab:res-robust} reporta a condição $f = 1{,}00$ sob as duas
sondas de deslocamento. Além da acurácia absoluta, listamos $R_c$ e a
\emph{habilidade retida}
\begin{equation}
  S_c = \frac{A^{c} - A^{\text{acaso}}}{A^{\text{clean}} - A^{\text{acaso}}},
  \label{eq:skill}
\end{equation}
que desconta o acerto trivial antes de formar a razão. A diferença entre as
duas colunas é instrutiva: $R_c$ dá crédito por acertar ao acaso, $S_c$ não.

\begin{table}[t]
  \centering
  \caption{Desempenho sob transformações não vistas em $f = 1{,}00$, sem
  augmentação em teste. $A^{c}$ é a acurácia absoluta (\%), $R_c$ a razão de
  retenção (\%) e $S_c$ a habilidade retida (\%) da Equação~\eqref{eq:skill}.
  Acaso: $10\%$ no CIFAR-10, $50\%$ no D-Fire.}
  \label{tab:res-robust}
  \small
  \begin{tabular}{llcccccc}
    \toprule
    \multirow{2}{*}{Modelo} & \multirow{2}{*}{Aug.}
      & \multicolumn{3}{c}{Rotação grande} & \multicolumn{3}{c}{Todas as transformações} \\
    \cmidrule(lr){3-5}\cmidrule(lr){6-8}
     & & $A^{c}$ & $R_c$ & $S_c$ & $A^{c}$ & $R_c$ & $S_c$ \\
    \midrule
    \multicolumn{8}{l}{\textit{CIFAR-10}} \\
    \multirow{3}{*}{Efficient-CapsNet}
      & nenhuma & 38,1 & 51,3 & 43,7 & 12,6 & 17,0 & 4,1 \\
      & padrão  & 52,9 & 62,4 & 57,4 & 12,8 & 15,1 & 3,8 \\
      & forte   & 51,4 & 61,9 & 56,6 & 13,8 & 16,7 & 5,2 \\
    \multirow{3}{*}{ResNet-18}
      & nenhuma & 36,4 & 46,2 & 38,4 & 11,6 & 14,7 & 2,3 \\
      & padrão  & 57,8 & 63,6 & 59,1 & 12,3 & 13,6 & 2,9 \\
      & forte   & \textbf{60,0} & 64,4 & 60,1 & 13,4 & 14,4 & 4,1 \\
    \multirow{3}{*}{DeiT-Tiny}
      & nenhuma & 40,0 & 65,8 & 59,0 & \textbf{16,4} & 26,9 & 12,6 \\
      & padrão  & 52,8 & 75,1 & 71,0 & 15,9 & 22,6 & 9,8 \\
      & forte   & 52,0 & \textbf{75,7} & \textbf{71,6} & \textbf{16,8} & \textbf{24,5} & \textbf{11,6} \\
    \midrule
    \multicolumn{8}{l}{\textit{D-Fire}} \\
    \multirow{3}{*}{Efficient-CapsNet}
      & nenhuma & 61,4 & 68,1 & 28,4 & 53,1 & 58,9 & 7,7 \\
      & padrão  & 77,6 & 82,9 & 63,2 & 56,1 & 59,9 & 13,9 \\
      & forte   & 80,0 & 87,0 & 71,5 & 60,5 & 65,8 & 25,0 \\
    \multirow{3}{*}{ResNet-18}
      & nenhuma & 63,4 & 69,5 & 32,4 & 52,5 & 57,5 & 6,0 \\
      & padrão  & 78,7 & 83,7 & 65,2 & 54,9 & 58,4 & 11,2 \\
      & forte   & \textbf{80,9} & 85,4 & 69,1 & \textbf{60,7} & 64,1 & 23,9 \\
    \multirow{3}{*}{DeiT-Tiny}
      & nenhuma & 63,5 & 75,0 & 38,9 & 52,9 & 62,5 & 8,4 \\
      & padrão  & 78,4 & 96,9 & 91,8 & 53,7 & 66,4 & 12,1 \\
      & forte   & 76,0 & \textbf{98,0} & \textbf{94,4} & 59,9 & \textbf{77,2} & \textbf{35,9} \\
    \bottomrule
  \end{tabular}
\end{table}

\subsection{A sonda de rotação}

A rotação grande é o teste direto de H1, e o resultado é desfavorável à
hipótese. No CIFAR-10 com augmentação padrão ou forte, a ResNet-18 tem a maior
acurácia absoluta sob rotação --- $57{,}8\%$ e $60{,}0\%$ contra $52{,}9\%$ e
$51{,}4\%$ do Efficient-CapsNet. A vantagem convolucional sob rotação
($+8{,}6$~pp em $f = 1{,}00$, augmentação forte) é maior, em valor absoluto, do
que qualquer vantagem capsular observada em toda a grade.

Apenas no regime sem augmentação nenhuma o Efficient-CapsNet fica à frente sob
rotação, e por pouco: $38{,}1\%$ contra $36{,}4\%$ da ResNet-18. Em teste
pareado por semente sobre todas as 27 células do CIFAR-10, a diferença de
retenção entre Efficient-CapsNet e ResNet-18 sob rotação é de $+1{,}1$~pp com
$p = 0{,}13$ --- indistinguível de zero. No D-Fire a mesma comparação dá
$-1{,}5$~pp com $p = 0{,}013$, ou seja, uma desvantagem capsular
estatisticamente detectável. Contra o DeiT-Tiny a diferença de retenção é de
$-15{,}1$~pp no CIFAR-10 e $-11{,}7$~pp no D-Fire, ambas com $p < 10^{-12}$.

A Figura~\ref{fig:res-dumbbell} torna o padrão visível: as barras do
Efficient-CapsNet e da ResNet-18 partem de pontos diferentes e caem
aproximadamente na mesma proporção, sem que a distância entre elas se feche.

\begin{figure}[t]
  \centering
  \includegraphics[width=\textwidth]{fig01_shift_dumbbell_cifar_10.png}
  \caption{CIFAR-10: acurácia limpa (círculo aberto) e acurácia sob
  deslocamento (marcador preenchido) para cada combinação de arquitetura,
  fração de dados e regime de augmentação. O comprimento da barra é a
  degradação absoluta.}
  \label{fig:res-dumbbell}
\end{figure}

\subsection{A composição de todas as transformações é um piso, não uma medida}
\label{sec:res:floor}

Sob a composição das seis corrupções, os três modelos deixam de funcionar. No
CIFAR-10 a acurácia cai para a faixa de $11{,}1\%$ a $19{,}1\%$, contra
$10{,}0\%$ de acaso, em todas as 54 células. No D-Fire os valores ficam entre
$-2{,}2$ e $+7{,}8$~pp da linha de base da classe majoritária ($53{,}4\%$); sem
augmentação de treino, 14 das 18 células ficam \emph{abaixo} dela, isto é,
perde para um classificador que ignora a imagem
(Figura~\ref{fig:res-baselines}).

Isso muda o que se pode concluir a partir dessa condição. Comparar retenção
entre modelos que estão todos no nível do acaso compara padrões de colapso, não
robustez. As diferenças de $R_c$ que aparecem ali --- $26{,}9\%$ para o
DeiT-Tiny contra $14{,}7\%$ para a ResNet-18, por exemplo --- correspondem, em
acurácia absoluta, a $16{,}4\%$ contra $11{,}6\%$: quatro pontos e meio acima e
um ponto e meio acima do acaso, respectivamente. Nenhum dos dois é utilizável.

\begin{figure}[t]
  \centering
  \includegraphics[width=\textwidth]{fig11_cross_dataset_vs_baseline.png}
  \caption{Os dois conjuntos contra suas próprias linhas de base triviais.
  Barras abaixo da linha pontilhada perdem para um classificador que ignora a
  imagem. Sob todas as transformações, nenhuma arquitetura se sustenta em
  nenhum dos conjuntos.}
  \label{fig:res-baselines}
\end{figure}

A recomendação que decorre daí é metodológica e vale além deste trabalho:
composições de corrupções só informam enquanto os modelos comparados
permanecerem acima do acaso por margem apreciável. Abaixo disso, a métrica
mede outra coisa. A Seção~\ref{sec:res:collapse} identifica exatamente o que.

\subsection{Por que a retenção não pode ser a métrica principal}
\label{sec:res:confound}

O DeiT-Tiny tem a maior razão de retenção em praticamente toda a
Tabela~\ref{tab:res-robust} e a menor acurácia limpa em toda a
Tabela~\ref{tab:res-clean}. Não é coincidência. A Figura~\ref{fig:res-confound}
mostra a correlação entre acurácia limpa e retenção no CIFAR-10:
$r = -0{,}44$ sob rotação e $r = -0{,}91$ sob todas as transformações. As
inclinações ajustadas são de $-34$ e $-48$~pp de retenção por unidade de
acurácia limpa.

\begin{figure}[t]
  \centering
  \includegraphics[width=\textwidth]{fig03_retention_confound_cifar_10.png}
  \caption{CIFAR-10: razão de retenção contra acurácia limpa. A relação
  negativa significa que a retenção, isolada, recompensa modelos que já eram
  ruins em distribuição.}
  \label{fig:res-confound}
\end{figure}

O mecanismo é aritmético. A acurácia sob deslocamento tem piso no acaso,
enquanto a acurácia limpa não tem teto próximo; um modelo que sai de $60\%$ e
cai para $40\%$ retém $67\%$, e um que sai de $93\%$ e cai para $60\%$ retém
$64\%$ --- embora o segundo esteja $20$~pp à frente onde importa. A
Figura~\ref{fig:res-frontier} mostra a mesma informação no plano
(acurácia limpa, acurácia deslocada): sob rotação a correlação entre os dois
eixos é $+0{,}68$, e sob todas as transformações ela \emph{inverte} de sinal
($-0{,}74$), que é a assinatura do regime de colapso descrito acima.

\begin{figure}[t]
  \centering
  \includegraphics[width=\textwidth]{fig02_robustness_frontier_cifar_10.png}
  \caption{CIFAR-10: acurácia sob deslocamento contra acurácia limpa. A
  diagonal tracejada marca invariância perfeita. Separação vertical entre cores
  na mesma abscissa seria o efeito arquitetural puro; ela é pequena sob rotação
  e some sob todas as transformações.}
  \label{fig:res-frontier}
\end{figure}

Por isso a Tabela~\ref{tab:res-robust} reporta $A^c$ em primeiro lugar, e por
isso H1 é avaliada sobre acurácia absoluta e sobre retenção pareada por
semente, não sobre médias de $R_c$.

% ---------------------------------------------------------------------
\section{Colapso de predição e detecção de erro}
\label{sec:res:collapse}

Dois diagnósticos separam ``degradação graciosa'' de ``falha silenciosa'': a
fração de predições que caem na classe mais predita (\emph{top-pred share}) e a
AUROC da confiança do modelo como detector dos próprios erros. A
Tabela~\ref{tab:res-auroc} traz os dois.

\begin{table}[t]
  \centering
  \caption{AUROC de detecção de erro (a confiança separa acertos de erros?) e
  fração de predições na classe mais predita, médias sobre as três frações de
  dados, sem augmentação em teste. AUROC de $0{,}50$ é o acaso; \emph{share} de
  $0{,}10$ (CIFAR-10) e $0{,}53$ (D-Fire) corresponde a predições
  balanceadas.}
  \label{tab:res-auroc}
  \footnotesize
  \setlength{\tabcolsep}{4pt}
  \begin{tabular}{lccccccccc}
    \toprule
    \multirow{2}{*}{Modelo}
      & \multicolumn{3}{c}{Limpo} & \multicolumn{3}{c}{Rotação grande} & \multicolumn{3}{c}{Todas} \\
    \cmidrule(lr){2-4}\cmidrule(lr){5-7}\cmidrule(lr){8-10}
     & nenh. & padr. & forte & nenh. & padr. & forte & nenh. & padr. & forte \\
    \midrule
    \multicolumn{10}{l}{\textit{CIFAR-10 --- AUROC}} \\
    Efficient-CapsNet & 0,799 & 0,856 & 0,848 & 0,707 & 0,762 & 0,766 & 0,533 & 0,547 & 0,548 \\
    ResNet-18         & 0,788 & 0,851 & 0,867 & 0,697 & 0,790 & 0,804 & 0,540 & 0,527 & 0,540 \\
    DeiT-Tiny         & 0,716 & 0,736 & 0,720 & 0,677 & 0,722 & 0,713 & 0,560 & 0,587 & 0,618 \\
    \addlinespace
    \multicolumn{10}{l}{\textit{CIFAR-10 --- fração na classe mais predita}} \\
    Efficient-CapsNet & 0,12 & 0,11 & 0,11 & 0,24 & 0,18 & 0,20 & 0,39 & 0,48 & 0,43 \\
    ResNet-18         & 0,13 & 0,11 & 0,11 & 0,28 & 0,18 & 0,18 & 0,55 & 0,70 & 0,76 \\
    DeiT-Tiny         & 0,13 & 0,12 & 0,13 & 0,21 & 0,15 & 0,16 & 0,38 & 0,48 & 0,38 \\
    \midrule
    \multicolumn{10}{l}{\textit{D-Fire --- AUROC}} \\
    Efficient-CapsNet & 0,837 & 0,838 & 0,836 & 0,581 & 0,696 & 0,728 & 0,515 & 0,558 & 0,595 \\
    ResNet-18         & 0,826 & 0,845 & 0,881 & 0,610 & 0,736 & 0,745 & 0,532 & 0,562 & 0,601 \\
    DeiT-Tiny         & 0,769 & 0,755 & 0,731 & 0,627 & 0,727 & 0,726 & 0,525 & 0,560 & 0,566 \\
    \addlinespace
    \multicolumn{10}{l}{\textit{D-Fire --- fração na classe mais predita}} \\
    Efficient-CapsNet & 0,56 & 0,54 & 0,55 & 0,78 & 0,64 & 0,62 & 0,75 & 0,60 & 0,58 \\
    ResNet-18         & 0,53 & 0,53 & 0,53 & 0,70 & 0,61 & 0,62 & 0,65 & 0,58 & 0,62 \\
    DeiT-Tiny         & 0,55 & 0,54 & 0,54 & 0,55 & 0,54 & 0,54 & 0,70 & 0,63 & 0,63 \\
    \bottomrule
  \end{tabular}
\end{table}

Em dados limpos, o Efficient-CapsNet tem a melhor detecção de erro do CIFAR-10
sob augmentação padrão ($0{,}856$ contra $0{,}851$ da ResNet-18) e a melhor do
D-Fire no regime sem augmentação ($0{,}837$ contra $0{,}826$), ficando a menos
de um centésimo da ResNet-18 sob augmentação padrão. A norma do vetor de cápsula, portanto,
funciona como escore de confiança pelo menos tão bem quanto o \emph{softmax},
o que é um resultado positivo para o modelo capsular e independente da
acurácia.

Sob deslocamento essa propriedade não sobrevive. Sob rotação a AUROC capsular
cai para $0{,}707$--$0{,}766$ no CIFAR-10 e para $0{,}581$ no D-Fire sem
augmentação --- neste último caso, com $77{,}6\%$ das predições em uma única
classe. Ou seja: o modelo rotacionado no D-Fire responde quase sempre a mesma
coisa e não sabe que está errando. Sob todas as transformações, a AUROC dos
três modelos fica entre $0{,}515$ e $0{,}618$, praticamente o acaso
(Figura~\ref{fig:res-auroc}). A falha, nessa condição, é silenciosa em todas as
arquiteturas.

\begin{figure}[t]
  \centering
  \includegraphics[width=\textwidth]{fig07_error_detection_auroc_cifar_10.png}
  \caption{CIFAR-10: AUROC da confiança como detector dos próprios erros. A
  linha vermelha marca o acaso. Sob a composição de transformações, nenhum
  modelo distingue acerto de erro.}
  \label{fig:res-auroc}
\end{figure}

A fração na classe mais predita revela a diferença de estilo de falha. Sob
todas as transformações no CIFAR-10, a ResNet-18 colapsa: $76\%$ das predições
em uma única classe no regime de augmentação forte, contra $43\%$ do
Efficient-CapsNet e $38\%$ do DeiT-Tiny (Figura~\ref{fig:res-collapse}). A
Figura~\ref{fig:res-perclass} mostra a forma desse colapso por classe: a
ResNet-18 corrompida acerta $42\%$ de \emph{airplane} e $55\%$ de \emph{bird} e
fica em $0$--$3\%$ em sete das dez classes. O Efficient-CapsNet degrada de
modo mais distribuído, sem por isso acertar mais.

\begin{figure}[t]
  \centering
  \includegraphics[width=\textwidth]{fig08_prediction_collapse_cifar_10.png}
  \caption{CIFAR-10: fração das predições concentradas na classe mais predita.
  Valores próximos de $1/K$ indicam predições ainda distribuídas; valores altos
  indicam colapso.}
  \label{fig:res-collapse}
\end{figure}

\begin{figure}[t]
  \centering
  \includegraphics[width=\textwidth]{fig09_per_class_accuracy_cifar_10.png}
  \caption{CIFAR-10: acurácia por classe e condição de teste, sem augmentação
  de treino, $f = 1{,}00$, agregada sobre sementes.}
  \label{fig:res-perclass}
\end{figure}

Um detalhe da Figura~\ref{fig:res-perclass} merece nota. A classe \emph{frog} é
a única que o Efficient-CapsNet mantém sob rotação grande: $85\%$ limpo e
$85\%$ rotacionado, retenção integral. A ResNet-18 mostra o mesmo padrão nessa
classe ($86\% \to 80\%$). A explicação mais econômica não é arquitetural, e sim
do dado: sapos aparecem em poses e orientações variadas e sem eixo vertical
dominante, ao contrário de cavalos, navios e caminhões. Onde o prior de
orientação é fraco, a rotação não penaliza --- em qualquer arquitetura.

% ---------------------------------------------------------------------
\section{Eficiência de dados}
\label{sec:res:data}

A Tabela~\ref{tab:res-dataeff} traz a estatística $\Delta_f$ da
Equação~\eqref{eq:data-eff}, a queda de acurácia limpa entre $f = 1{,}00$ e
$f = 0{,}33$.

\begin{table}[t]
  \centering
  \caption{Eficiência de dados: $\Delta_f = A_{f=1{,}00} - A_{f=0{,}33}$ em pp
  de acurácia limpa. Valores menores indicam menor dependência do volume de
  dados.}
  \label{tab:res-dataeff}
  \small
  \begin{tabular}{llccc c}
    \toprule
    \multirow{2}{*}{Conjunto} & \multirow{2}{*}{Modelo} & \multicolumn{3}{c}{Augmentação} & \multirow{2}{*}{Média} \\
    \cmidrule(lr){3-5}
     & & nenhuma & padrão & forte & \\
    \midrule
    \multirow{3}{*}{CIFAR-10}
      & Efficient-CapsNet & 10,4 & \phantom{0}6,7 & \phantom{0}6,5 & \textbf{7,9} \\
      & ResNet-18         & 12,6 & \phantom{0}7,5 & \phantom{0}6,5 & 8,9 \\
      & DeiT-Tiny         & \phantom{0}8,1 & 11,0 & 19,0 & 12,7 \\
    \midrule
    \multirow{3}{*}{D-Fire}
      & Efficient-CapsNet & \phantom{0}2,9 & \phantom{0}4,6 & \phantom{0}5,4 & 4,3 \\
      & ResNet-18         & \phantom{0}6,8 & \phantom{0}5,3 & \phantom{0}4,7 & 5,6 \\
      & DeiT-Tiny         & \phantom{0}2,4 & \phantom{0}5,0 & \phantom{0}3,6 & \textbf{3,7} \\
    \bottomrule
  \end{tabular}
\end{table}

O Efficient-CapsNet tem o menor $\Delta_f$ médio no CIFAR-10 ($7{,}9$ contra
$8{,}9$~pp da ResNet-18) e fica $1{,}3$~pp à frente dela no D-Fire. A direção é
a prevista por H2, mas a magnitude não sustenta a hipótese: no teste pareado
por semente e regime de augmentação, a diferença Efficient-CapsNet menos
ResNet-18 é de $-1{,}0$~pp com $p = 0{,}080$ no CIFAR-10 e $-1{,}3$~pp com
$p = 0{,}139$ no D-Fire, ambas com $n = 9$. Nenhuma alcança significância
convencional, e o poder com três sementes é baixo demais para que a ausência de
significância seja informativa em qualquer direção.

Além disso, o ordenamento absoluto não se altera em nenhuma fração: a
ResNet-18 em $f = 0{,}33$ já supera o Efficient-CapsNet em $f = 1{,}00$ no
CIFAR-10 sob augmentação forte ($86{,}8\%$ contra $83{,}0\%$). Uma eventual
inclinação mais suave da curva capsular não se converte em vantagem prática
dentro da faixa investigada.

A Figura~\ref{fig:res-scaling} mostra também um efeito que $\Delta_f$ não
captura: sob rotação grande, mais dados continuam ajudando a ResNet-18
($+4{,}9$~pp de $f = 0{,}33$ a $f = 1{,}00$ com augmentação forte), enquanto
sob todas as transformações as curvas são planas ou levemente decrescentes para
os três modelos. Dados adicionais compram acurácia em distribuição e sob
deslocamento moderado; não compram nada no regime de colapso.

\begin{figure}[t]
  \centering
  \includegraphics[width=\textwidth]{fig05_data_scaling_cifar_10.png}
  \caption{CIFAR-10: acurácia contra fração de dados de treino, por regime de
  augmentação (linhas) e condição de teste (colunas). Barras de erro são o
  desvio padrão entre três sementes.}
  \label{fig:res-scaling}
\end{figure}

% ---------------------------------------------------------------------
\section{Interação com a augmentação}
\label{sec:res:aug}

H3 previa que uma eventual vantagem capsular encolheria ao nível do ruído
quando as linhas de base recebessem augmentação forte. O resultado observado é
mais forte que a hipótese e aponta na direção contrária: não há vantagem
capsular a encolher, e a \emph{desvantagem} capsular \emph{cresce} com a
augmentação.

A Figura~\ref{fig:res-ladder} quantifica o quanto cada arquitetura converte
augmentação em desempenho. Subindo de ``nenhuma'' a ``forte'' no CIFAR-10, a
ResNet-18 ganha $17$~pp em dados limpos e $22$~pp sob rotação; o
Efficient-CapsNet ganha $11$ e $12$~pp; o DeiT-Tiny ganha $1$ e $9$~pp. A
ResNet-18 é, das três, a que melhor aproveita o regime.

\begin{figure}[t]
  \centering
  \includegraphics[width=\textwidth]{fig04_augmentation_ladder_cifar_10.png}
  \caption{CIFAR-10: efeito da intensidade da augmentação de treino sobre cada
  condição de teste, com média sobre as frações de dados. As anotações dão o
  ganho total de ``nenhuma'' a ``forte''.}
  \label{fig:res-ladder}
\end{figure}

A consequência aparece na Figura~\ref{fig:res-heatmap}, que mostra a diferença
entre o Efficient-CapsNet e a melhor linha de base em cada uma das 54
condições do CIFAR-10. Sem augmentação e com $f = 0{,}33$, a diferença é de
$-2{,}3$~pp em dados limpos e chega a $+2{,}3$~pp sob rotação, ambas dentro do
dobro do ruído entre sementes. Com augmentação forte e $f = 1{,}00$, ela é de
$-10{,}3$~pp em dados limpos e $-8{,}7$~pp sob rotação. Em média sobre as
frações, a desvantagem capsular em dados limpos passa de $-2{,}5$~pp (sem
augmentação) para $-8{,}8$ (padrão) e $-13{,}2$~pp (forte).

\begin{figure}[t]
  \centering
  \includegraphics[width=\textwidth]{fig06_advantage_heatmap_cifar_10.png}
  \caption{CIFAR-10: diferença de acurácia entre o Efficient-CapsNet e a melhor
  linha de base em cada condição. Azul indica vantagem capsular, vermelho
  desvantagem; o símbolo $^\circ$ marca células dentro do dobro do ruído entre
  sementes.}
  \label{fig:res-heatmap}
\end{figure}

A interpretação natural é que a invariância que as cápsulas ofereceriam por
construção pode ser adquirida pelas linhas de base por meio de dados, e que
adquiri-la por dados é mais barato e mais eficaz. O resultado é consistente com
a leitura crítica da literatura: o que se credita ao roteamento talvez seja, em
boa parte, atribuível ao objetivo de treino e ao regime de augmentação
(Seção~\ref{sec:protocol:threats}).

Vale registrar o contraponto: nem o Efficient-CapsNet nem o DeiT-Tiny se
beneficiam de augmentação forte em dados limpos; ambos pioram em relação à
augmentação padrão, nos dois conjuntos. Isso
sugere que a perda de margem com reconstrução, ou a geometria das cápsulas
primárias, tolera mal o RandAugment e o Random Erasing --- e que a comparação
sob ``augmentação forte'' talvez não esteja comparando cada arquitetura em seu
melhor. É uma limitação do delineamento, não um resultado sobre cápsulas.

% ---------------------------------------------------------------------
\section{Variabilidade entre sementes}
\label{sec:res:seeds}

A Tabela~\ref{tab:res-seeds} reporta a amplitude entre as três sementes, média
sobre todas as condições de avaliação.

\begin{table}[t]
  \centering
  \caption{Amplitude (máximo menos mínimo) da acurácia entre três sementes, em
  pp, média sobre todas as condições de avaliação.}
  \label{tab:res-seeds}
  \small
  \begin{tabular}{llcccc}
    \toprule
    \multirow{2}{*}{Conjunto} & \multirow{2}{*}{Modelo} & \multicolumn{3}{c}{Augmentação} & \multirow{2}{*}{Média} \\
    \cmidrule(lr){3-5}
     & & nenhuma & padrão & forte & \\
    \midrule
    \multirow{3}{*}{CIFAR-10}
      & Efficient-CapsNet & 0,93 & 0,93 & 0,86 & \textbf{0,91} \\
      & ResNet-18         & 2,01 & 1,17 & 1,06 & 1,42 \\
      & DeiT-Tiny         & 1,72 & 3,39 & 6,75 & 3,95 \\
    \midrule
    \multirow{3}{*}{D-Fire}
      & Efficient-CapsNet & 1,44 & 1,62 & 1,38 & \textbf{1,48} \\
      & ResNet-18         & 1,84 & 1,26 & 1,79 & 1,63 \\
      & DeiT-Tiny         & 2,14 & 7,57 & 7,70 & 5,80 \\
    \bottomrule
  \end{tabular}
\end{table}

Este é o resultado mais nítido a favor do modelo capsular em toda a grade, e
não é sobre robustez nem sobre acurácia. O Efficient-CapsNet é a arquitetura
mais estável entre sementes nos dois conjuntos, com amplitude média de
$0{,}91$~pp no CIFAR-10 --- menos de dois terços da ResNet-18 e menos de um
quarto do DeiT-Tiny. A estabilidade é quase insensível ao regime de
augmentação, enquanto a do DeiT-Tiny se deteriora fortemente com ele (de
$1{,}7$ para $6{,}8$~pp no CIFAR-10).

A amplitude individual mais extrema da grade é de $19{,}4$~pp, no DeiT-Tiny com
augmentação forte e $f = 1{,}00$ no CIFAR-10 --- uma única semente separa
$57\%$ de $76\%$ de acurácia. Os doze piores casos de dispersão da
grade pertencem todos ao DeiT-Tiny, oito deles no D-Fire. Qualquer número reportado para essa arquitetura neste
trabalho deve ser lido com essa dispersão em mente, e a
Figura~\ref{fig:res-seeds} torna explícito que várias comparações envolvendo
o DeiT-Tiny não sobrevivem ao ruído.

\begin{figure}[t]
  \centering
  \includegraphics[width=\textwidth]{fig10_seed_consistency_cifar_10.png}
  \caption{CIFAR-10, $f = 1{,}00$: cada semente desenhada individualmente,
  ligando a acurácia limpa à acurácia sob todas as transformações.}
  \label{fig:res-seeds}
\end{figure}

Uma leitura possível é que a perda de margem com reconstrução, por regularizar
mais fortemente, produz um objetivo com paisagem menos sensível à
inicialização. A hipótese não é testável dentro deste delineamento, que
confunde arquitetura e função objetivo, mas é a explicação candidata a ser
investigada.

% ---------------------------------------------------------------------
\section{Contagem de vitórias}
\label{sec:res:wins}

Somando as 108 condições de avaliação (2 conjuntos $\times$ 3 augmentações
$\times$ 3 frações $\times$ 2 augmentações em teste $\times$ 3 transformações),
a ResNet-18 vence em 57, o DeiT-Tiny em 33 e o Efficient-CapsNet em 18. Das 18
vitórias capsulares, \emph{uma única} excede o dobro do ruído entre sementes:
D-Fire, augmentação forte, $f = 1{,}00$, teste com augmentação forte e todas as
transformações, $59{,}2\%$ contra $57{,}7\%$ da ResNet-18 --- e essa é
justamente uma condição em que ambos os modelos estão a poucos pontos da linha
de base da classe majoritária.

A distribuição das vitórias também localiza os regimes. As 33 vitórias do
DeiT-Tiny concentram-se sob deslocamento (18 delas na composição de todas as
transformações no CIFAR-10, onde todos estão perto do acaso). As 57 vitórias da
ResNet-18 concentram-se em dados limpos e sob rotação, isto é, nos regimes em
que os números significam alguma coisa.

% ---------------------------------------------------------------------
\section{Veredito sobre as hipóteses}
\label{sec:res:verdict}

\begin{table}[t]
  \centering
  \caption{Resumo da avaliação das hipóteses da
  Seção~\ref{sec:protocol:hypotheses}.}
  \label{tab:res-verdict}
  \small
  \begin{tabular}{l p{0.68\textwidth}}
    \toprule
    Hipótese & Veredito \\
    \midrule
    H1 & \textbf{Rejeitada.} Sob rotação grande, a retenção capsular não supera
    a da ResNet-18 no CIFAR-10 ($+1{,}1$~pp, $p = 0{,}13$) e é inferior no
    D-Fire ($-1{,}5$~pp, $p = 0{,}013$). Em acurácia absoluta, a ResNet-18 está
    à frente em seis das nove condições do CIFAR-10, incluindo as seis com
    augmentação. \\
    \addlinespace
    H2 & \textbf{Não sustentada.} A direção prevista aparece ($\Delta_f$ menor
    para as cápsulas em ambos os conjuntos), mas a diferença é de cerca de
    $1$~pp e não atinge significância com três sementes
    ($p = 0{,}080$ e $p = 0{,}139$). O ordenamento absoluto não se inverte em
    nenhuma fração. \\
    \addlinespace
    H3 & \textbf{Rejeitada, com inversão de sinal.} Não há vantagem capsular a
    encolher. A desvantagem em dados limpos no CIFAR-10 \emph{cresce} de
    $-2{,}5$~pp sem augmentação para $-13{,}2$~pp sob augmentação forte. \\
    \bottomrule
  \end{tabular}
\end{table}

O achado central é negativo e deve ser enunciado sem atenuação: nesta grade,
substituir camadas convolucionais por cápsulas com roteamento por
autoatenção não compra robustez a transformações não vistas, nem eficiência de
dados suficiente para compensar a perda de acurácia em distribuição. Onde o
modelo capsular aparece à frente, ou a margem está dentro do ruído entre
sementes, ou os modelos comparados estão no nível do acaso.

Três achados positivos sobrevivem à mesma análise, e nenhum deles é sobre
equivariância:

\begin{description}
\item[Estabilidade.] O Efficient-CapsNet é a arquitetura menos sensível à
semente nos dois conjuntos, com amplitude média de $0{,}91$~pp no CIFAR-10
contra $1{,}42$ da ResNet-18 e $3{,}95$ do DeiT-Tiny
(Seção~\ref{sec:res:seeds}).

\item[Confiança calibrada em distribuição.] A norma do vetor de cápsula é um
escore de detecção de erro competitivo ou superior ao \emph{softmax} em dados
limpos, com AUROC de até $0{,}856$ no CIFAR-10 e $0{,}838$ no D-Fire
(Seção~\ref{sec:res:collapse}).

\item[Competitividade na tarefa binária.] No D-Fire, a diferença para a melhor
linha de base é de $0{,}8$--$3{,}4$~pp conforme o regime, contra
$2{,}5$--$13{,}2$~pp no CIFAR-10. A penalidade capsular é fortemente dependente
da tarefa, o que sugere que conclusões obtidas apenas no CIFAR-10 não
transferem.
\end{description}

% ---------------------------------------------------------------------
\section{Ameaças à validade destes resultados}
\label{sec:res:threats}

Além das ameaças já declaradas na Seção~\ref{sec:protocol:threats}, a análise
deste capítulo acrescenta quatro.

\begin{description}

\item[A composição de corrupções satura.] Sob todas as transformações, os três
modelos ficam no nível do acaso (CIFAR-10) ou da classe majoritária (D-Fire).
Comparações nessa condição não separam robustez; separam estilo de colapso.
Trinta e seis das 108 condições da Seção~\ref{sec:res:wins} pertencem a esse
regime e foram desconsideradas nas conclusões substantivas.

\item[Uma única magnitude de rotação.] A sonda amostra
$\theta \sim \mathcal{U}(-90^{\circ}, 90^{\circ})$ uma vez por imagem. Uma
varredura determinística em $\theta$ produziria uma curva, e é possível que as
cápsulas se saiam melhor em ângulos pequenos e intermediários, onde a
degradação convolucional ainda é suave. O resultado aqui reportado vale para o
regime agregado, não para cada ângulo.

\item[Poder estatístico.] Com três sementes, diferenças abaixo de
aproximadamente $2$~pp são indistinguíveis do ruído na maior parte da grade.
Isso protege contra falsos positivos e impede, simetricamente, afirmar que H2
é falsa: só é possível afirmar que ela não foi demonstrada.

\item[A augmentação forte não é neutra.] Das três arquiteturas, só a ResNet-18
tem a augmentação forte como melhor regime; o Efficient-CapsNet tem a padrão,
nos dois conjuntos. Comparar as três sob augmentação forte é comparar a
ResNet-18 em seu melhor com o Efficient-CapsNet fora dele. A comparação mais favorável ao modelo capsular --- melhor regime
contra melhor regime --- ainda o deixa $8{,}5$~pp atrás no CIFAR-10 e
$1{,}1$~pp atrás no D-Fire, de modo que a conclusão não se inverte; mas as
magnitudes da Seção~\ref{sec:res:aug} são um teto, não uma estimativa pontual.

\end{description}
